\exercisesfor{2}

\begin{enumerate}
\def\labelenumi{\arabic{enumi}.}
\tightlist
\item
  设 g 是具有基本族 \((x_{1},\ldots,x_{n},\ldots)\) 的自由李代数。记 \(g_{n}\) 为由 \(x_{1},\ldots,x_{n}\) 生成的 g 的子代数，并记 \(b_{n}\) 为 \(g_{n}\) 中包含 \(x_{n}\) 的最小理想。
\end{enumerate}

\begin{enumerate}
\def\labelenumi{(\alph{enumi})}
\tightlist
\item
  证明 \(\mathfrak{b}_n\) 是由下列元素生成的 \(\mathfrak{g}_n\) 的子模：
\end{enumerate}

\[
\left(\operatorname{ad} \left(x _ {i _ {1}}\right) \circ \dots \circ \operatorname{ad} \left(x _ {i _ {k}}\right)\right) \left(x _ {n}\right),
\]

其中 \(k \geqslant 0\)，且 \(i_h \leqslant n\) 对所有 \(h\) 都成立。

\begin{enumerate}
\def\labelenumi{(\alph{enumi})}
\setcounter{enumi}{1}
\item
  证明 \(\mathfrak{g}_n = \mathfrak{g}_{n-1} \oplus \mathfrak{b}_n\)；从而推出 \(\mathfrak{g}\) 是各个 \(\mathfrak{b}_n\)（ \(n \geqslant 1\)）的直和。
\item
  利用 (a) 和 (b)，证明 K-模 g 由元素 \([x_{i_{1}}, [x_{i_{2}}, \ldots, [x_{i_{k-1}}, x_{i_{k}}] \ldots]]\) 生成，其中 \(k \geqslant 0\)，且 \(i_{h} \leqslant i_{k}\) 对 h \textless{} k 成立。
\end{enumerate}

¶ 2. 设 X 是至少含有两个元素的可数集，H 是 M(X) 的 Hall 集子集所组成的集合（参见~定义 2）。证明 Card(H) = \(2^{N_0}\)。

\begin{enumerate}
\def\labelenumi{\arabic{enumi}.}
\setcounter{enumi}{2}
\item
  设 \(\mathbf{X}\) 是一个带全序的集合。证明存在一个相对于 \(\mathbf{X}\) 的 Hall 集，使得 \(\mathrm{H} \cap \mathrm{M}^3(\mathbf{X})\) 由乘积 \(z(yx)\)（其中 \(y < x\)、\(y \leqslant z\)）组成，并且 \(\mathrm{H} \cap \mathrm{M}^4(\mathbf{X})\) 由乘积 \(w(z(yx))\)（其中 \(w \geqslant z \geqslant y\)、\(y < x\)）以及乘积 \((ab)(cd)\)（其中 \(a < b\)、\(c < d\)，并且或者 \(a < c\)，或者 \(a = c\) 且 \(b < d\)）组成。
\item
  \begin{enumerate}
  \def\labelenumii{(\alph{enumii})}
  \tightlist
  \item
    证明由表示
  \end{enumerate}
\end{enumerate}

\[
\{x, y; [ x, [ x, y ] ] = [ y, [ x, y ] ] = 0 \}
\]

定义的李代数具有基 \((x,y,[x,y])\)。定义此代数到一个矩阵代数中的嵌入。

\begin{enumerate}
\def\labelenumi{(\alph{enumi})}
\setcounter{enumi}{1}
\tightlist
\item
  对由表示
\end{enumerate}

\[
\{x, y; [ x, [ x, y ] ] - 2 x = [ y, [ x, y ] ] + 2 y = 0 \}.
\]

\begin{enumerate}
\def\labelenumi{(\alph{enumi})}
\setcounter{enumi}{2}
\tightlist
\item
  证明具有表示
\end{enumerate}

\[
\{x, y, z; [ x, y ] - x = [ y, z ] - y = [ z, x ] - z = 0 \}
\]

的李代数化为 0。

\begin{enumerate}
\def\labelenumi{\arabic{enumi}.}
\setcounter{enumi}{4}
\item
  设 \(\mathfrak{g}\) 是自由李代数，\(\mathfrak{r}\) 是 \(\mathfrak{g}\) 的理想。假设 \(\mathfrak{r} = [\mathfrak{g},\mathfrak{r}]\)。证明 \(\mathfrak{r} = \{0\}\)（利用 \(\bigcap_{n\geqslant 0}\mathcal{C}^n\mathfrak{g} = \{0\}\) 这一事实）。
\item
  设 \(\mathbf{E}\) 是一个模。令 \(\mathbf{ME}\) 为 \(\mathbf{E}\) 的自由代数（代数，第三章，§2，习题 13）。回忆 \(\mathrm{ME} = \bigoplus_{n\geqslant 1}\mathrm{E}_n\)，其中
\end{enumerate}

\[
\mathrm{E} _ {1} = \mathrm{E}, \qquad \mathrm{E} _ {n} = \bigoplus_ {p + q = n} \mathrm{E} _ {p} \otimes \mathrm{E} _ {q} \quad \text {if} \quad n \geqslant 2,
\]

且 ME 上的代数结构由典范映射 \(\mathbf{E}_n\otimes \mathbf{E}_m\to \mathbf{E}_{n + m}\) 定义。

\begin{enumerate}
\def\labelenumi{(\alph{enumi})}
\item
  设 J 是 ME 中包含形如 xx 以及 \((xy)z + (yz)x + (zx)y\)（其中 x、y、z 属于 E）的元素的最小双边理想。令 LE = ME/J。证明 J 是 ME 的分次理想；从而得到 LE 的一个分次 \((\mathrm{L}^{n}\mathrm{E})_{n\geqslant1}\)。
\item
  证明 LE 是一个李代数；它称为模 E 的自由李代数。证明对每个李代数 g 和每个线性映射 f: \(E \to g\)，存在唯一一个延拓 f 的李代数同态 F: \(LE \to g\)。
\item
  定义同构 \(\mathbf{E} \to \mathbf{L}^1\mathbf{E}\) 和 \(\wedge^2\mathbf{E} \to \mathbf{L}^2\mathbf{E}\)。证明 \(\mathbf{L}^3\mathbf{E}\) 可与 \(\mathbf{E} \otimes \wedge^2\mathbf{E}\) 对由下列元素生成的子模取商所得的模相识别：
\end{enumerate}

\[
x \otimes (y \wedge z) + y \otimes (z \wedge x) + z \otimes (x \wedge y) \quad \text { for } x, y, z \text { in } E.
\]

\begin{enumerate}
\def\labelenumi{(\alph{enumi})}
\setcounter{enumi}{3}
\item
  证明当 E 是以 X 为基的自由模时，LE 可与第 2 号中定义的自由李代数 \(\mathrm{L}(\mathbf{X})\) 相识别。
\item
  证明 E 到 LE 的包络代数 U(LE) 中的典范映射可以延拓为张量代数 TE 到 U(LE) 的一个同构。
\item
  设 \(\sigma\) 是从 \(\wedge^2\mathrm{E}\) 到 \(\mathsf{T}^2\mathsf{E}\) 的线性映射，满足
\end{enumerate}

\[
\sigma (x \wedge y) = x \otimes y - y \otimes x.
\]

构造一个使 \(\sigma\) 不是单射的模 E。由此推出，对于这个模，LE 到 U(LE) 的典范映射不是单射（参见第一章 §2 的习题 9）。

\begin{enumerate}
\def\labelenumi{\arabic{enumi}.}
\setcounter{enumi}{6}
\tightlist
\item
  设 K 是一个域。令 l 是具有基本族 \((x_{i})_{i\in I}\) 的自由李代数，并令 M 是一个 l-模。
\end{enumerate}

\begin{enumerate}
\def\labelenumi{(\alph{enumi})}
\item
  证明 \(\mathbf{H}^2 (\mathbf{l},\mathbf{M}) = \{0\}\)（参见~第一章 §3 的习题 12。使用命题 1 的推论 2 以及所述习题的 (i) 部分）。
\item
  证明对于 M 中元素的每个族 \((m_{i})_{i\in I}\)，存在唯一的一个 1 次协循环 \(\phi:l\to M\)，满足 \(\phi (x_i) = m_i\)。推出正合列：
\end{enumerate}

\[
0 \rightarrow \mathrm{H} ^ {0} (\mathfrak {l}, \mathrm{M}) \rightarrow \mathrm{M} \rightarrow \mathrm{M} ^ {\mathrm{I}} \rightarrow \mathrm{H} ^ {1} (\mathfrak {l}, \mathrm{M}) \rightarrow 0.
\]

若 I 的基数为有限数 \(n\)，且 M 在 K 上具有有限秩，则

\[
\operatorname{rg} \mathrm{H} ^ {1} (\mathfrak {l}, \mathrm{M}) - \operatorname{rg} \mathrm{H} ^ {0} (\mathfrak {l}, \mathrm{M}) = (n - 1) \operatorname{rg} \mathrm{M}.
\]

\begin{enumerate}
\def\labelenumi{\arabic{enumi}.}
\setcounter{enumi}{7}
\tightlist
\item
  设 K 是一个域。令 l 是具有基本族 \((\overline{x}_{i})_{i\in I}\) 的自由李代数，令 r 是包含于 {[}l, l{]} 中的 l 的理想，并令 g = l/r；用 \(x_{i}\) 表示 \(\overline{x}_{i}\) 在 g 中的像。
\end{enumerate}

证明以下性质等价：

\begin{enumerate}
\def\labelenumi{(\roman{enumi})}
\item
  g 是自由李代数。
\item
  g 以 \((x_{i})\) 为基本族（即~r = \{0\}）。
\item
  对每个 g-模 M，都有 H²(g, M) = \{0\}。
\item
  若 g 在 K 上平凡作用，则 \(H^{2}(g, K) = \{0\}\)。
\end{enumerate}

（蕴含关系 (ii) \(\Rightarrow\) (i) 和 (iii) \(\Rightarrow\) (iv) 显然，而 (i) \(\Rightarrow\) (iii) 由习题 7 得出。为证明 (iv) \(\Rightarrow\) (ii)，只需证明 \(\mathfrak{r} = [\mathfrak{l},\mathfrak{r}]\)，参见~习题 5；否则取 \(\mathfrak{h}\) 为 \(\mathfrak{r}\) 的一个包含 \([\mathfrak{l},\mathfrak{r}]\) 的超平面，并注意到扩张 \(\mathfrak{r} / \mathfrak{h}\to \mathfrak{l} / \mathfrak{h}\to \mathfrak{l} / \mathfrak{r} = \mathfrak{g}\) 是本质的。） 

¶ 9. 设 \(\mathfrak{g} = \bigoplus_{n \geqslant 1} \mathfrak{g}_n\) 是一个分次李代数。若 \(\mathfrak{h}\) 是 \(\mathfrak{g}\) 的分次子代数，且 \(\mathfrak{g} = \mathfrak{h} + [\mathfrak{g}, \mathfrak{g}]\)，证明 \(\mathfrak{h} = \mathfrak{g}\)。若 \((x_i)\) 是 \(\mathfrak{g}\) 的齐次元素族，则推出 \((x_i)\) 是生成族，当且仅当 \(x_i\) 在 \(\mathfrak{g}/[\mathfrak{g}, \mathfrak{g}]\) 中的像生成 \(\mathfrak{g}/[\mathfrak{g}, \mathfrak{g}]\)。

若 K 是 g 在其上平凡作用的域，且 \(\mathrm{H}^{2}(\mathfrak{g},\mathbf{K})=\{0\}\)，证明族 \((x_{i})\) 是基本族，当且仅当 \(x_{i}\) 在 g/{[}g, g{]} 中的像构成这个向量空间的一组基。（使用习题 8。）†

\begin{enumerate}
\def\labelenumi{\arabic{enumi}.}
\setcounter{enumi}{9}
\tightlist
\item
  设 \(t\) 是具有一个含 \(n\) 个元素的有限基本族的自由李代数，并令 \(u\) 是 \(l\) 的满射自同态。证明 \(u\) 是同构。（写出 \(\mathrm{gr}^n(l) = \mathcal{C}^n / \mathcal{C}^{n+1}l\)，并以 \(\mathrm{gr}^n(u)\) 表示 \(\mathrm{gr}^n(l)\) 中由 \(u\) 导出的自同态；注意 \(\mathrm{gr}^n(u)\) 是满射；由于 \(\mathrm{gr}^n(l)\) 是有限秩自由 K-模，推出 \(\mathrm{gr}^n(u)\) 是双射，因为 \(u\) 的核包含于 \(\bigcap_{n} \mathcal{C}^n l\)，而该交为 \(\{0\}\)。）
\end{enumerate}

若 \((y_{1},\ldots ,y_{m})\) 是 l 的生成族，证明 \(m\geqslant n\)，并证明等号成立当且仅当 \((y_{1},\ldots ,y_{m})\) 是基本族。

¶ 11. 设 Y 和 Z 是两个不相交的集合（Y 带有全序），令 \(X = Y \cup Z\) 且 \(l = L(X)\)。令 \(a_{Y}\) 为由 Z 和 {[}l, l{]} 生成的 l 的子模；它是李代数 l 的一个理想。我们拟显式求出 \(a_{Y}\) 的一个基本族。

令 \(\mathbf{M}\) 为 \(\mathbf{Y}\) 上有限支撑的正整数值函数的集合。若 \(m\in \mathbf{M}\)，用 \(\theta^m\) 表示由下式给出的 l 的自同态：

\[
\theta^ {m} = \prod_ {y \in Y} (a d y) ^ {m (y)},
\]

其中乘积按 Y 上给定的序关系进行。令 \(\mathfrak{S}\) 表示 \(\mathbf{Y} \times \mathbf{M}\) 的子集，它由满足下述条件的有序对 \((u, m)\) 构成：存在 \(y \in \mathbf{Y}\) 使得 \(y > u\) 且 \(m(y) \geqslant 1\)。证明元素

\[
\theta^ {m} (u) \quad \text { for } \quad (u, m) \in \mathfrak {S} \quad \text { and } \quad \theta^ {m} (z) \quad \text { for } \quad (z, m) \in Z \times M
\]

构成李代数 \(\mathfrak{a}_{\mathbf{Y}}\) 的一个基本族。

（将问题归结为 Y 有限的情形，并按 Card(Y) 归纳。将命题 10 的推论应用于 Y 的最大元 y，并将归纳假设应用于 \(Y = \{y\}\)。）

\begin{enumerate}
\def\labelenumi{\arabic{enumi}.}
\setcounter{enumi}{11}
\tightlist
\item
  设 \(\mathbf{X} = \{x, y\}\) 是一个含两个元素的集合。证明 \(\mathbf{L}(\mathbf{X})\) 的导出代数是自由李代数，并以元素 \(((\operatorname{ad}y)^q \circ (\operatorname{ad}x)^p)(y)\) 为基本族，其中 \(p \geqslant 1\)、\(q \geqslant 0\)。（使用习题 11。）
\end{enumerate}

¶ 13.（本习题中假设 K 上的每个投射模都是自由的；例如当 K 为主理想整环时即如此。）

令 \(\mathfrak{l} = \sum_{n=1}^{\infty} \mathfrak{l}_n\) 是一个具有由齐次元素组成的基本族 B 的分次李代数，并令 \(\mathfrak{h}\) 是 \(\mathfrak{l}\) 的一个分次李子代数；作为模看，它是 \(\mathfrak{l}\) 的一个直因子。

\begin{enumerate}
\def\labelenumi{(\alph{enumi})}
\tightlist
\item
  对所有 \(i \geqslant 0\)，令 \(\mathfrak{l}^{(i)}\) 为 \(\mathfrak{l}\) 的分次子代数，满足
\end{enumerate}

\[
\begin{array}{l l} \mathfrak {l} _ {j} ^ {(i)} = \mathfrak {h} _ {j} & \text { if } j \leqslant i. \\ \mathfrak {l} _ {j} ^ {(i)} = \mathfrak {l} _ {j} & \text { if } j > i. \end{array}
\]

于是 \(l = \ell^{(0)} \supset \ell^{(1)} \supset \cdots \supset \mathfrak{h}\)。证明通过对 \(i\) 归纳，存在一个记作 \(B^{(i)}\) 的 \(\ell^{(i)}\) 的基本族，且由齐次元素组成，使得 \(B^{(i-1)}\) 中以及 \(B^{(i)}\) 中次数 \(< i - 1\) 的元素相同。（假设 \(B^{(i-1)}\) 已构造。令 \(m_i\) 为 \(l_i = l_i^{(i-2)}\) 与由 \(b_j\)（其中 \(j < i\)）生成的子代数的交，并令 \(b_i\) 为 \(l_i\) 中由 \(B^{(i-1)}\) 中次数为 \(i\) 的元素生成的子模。归纳假设蕴含 \(l_i = m_i \oplus b_i\)。由于 \(m_i \subset b_i\)，可以将 \(b_i\) 分解为 \(b_i = y_i \oplus b_i\)，从而 \(b_i = m_i \oplus b_i\)；由关于 K 的假设，模 \(y_i\) 和 \(b_i\) 是自由的；

必要时改变 \(\mathrm{B}^{(i-1)}\) 中次数为 i 的元素，可以假设 \(b_{i}\) 由 \(\mathrm{B}^{(i-1)}\) 的一个子集生成。将习题 11 应用于 \(\mathfrak{l}^{(i-1)}\) 及其理想 \(\mathfrak{l}^{(i)}\)，得到一个记作 \(\mathrm{B}^{(i)}\) 的 \(\mathfrak{l}^{(i)}\) 的基本族，具有所需性质。）

\begin{enumerate}
\def\labelenumi{(\alph{enumi})}
\setcounter{enumi}{1}
\tightlist
\item
  由 (a) 推出 \(\mathfrak{h}\) 具有由齐次元素组成的基本族。
\end{enumerate}

¶ 14. 设 K 是一个域。令 X 是一个集合，x、y 是 L(X) 中关于 K 线性无关的两个元素。证明族 \((x,y)\) 在李代数 L(X) 中是自由的。（令 \(x_{p}\)（相应地 \(y_{q}\)）为 x（相应地 y）的最高次数非零齐次分量。必要时将 y 的一个倍数加到 x 上，可以假设 \(x_{p}\) 和 \(y_{q}\) 线性无关。由 \(x_{p}\) 和 \(y_{q}\) 生成的 L(X) 的子代数是分次的，因而是自由的，参见~习题 13。推出 \((x_{p},y_{q})\) 是自由族，再由此过渡到 \((x,y)\)。）

¶ 15. 设 \(\mathrm{X} = \{x, y\}\) 是一个含两个元素的集合，且 \(\sigma\) 是李代数 \(\mathrm{L}(\mathrm{X})\) 的一个自同构。证明若 \(\mathrm{K}\) 是域，则 \(\sigma\) 保持 \(\mathrm{L}(\mathrm{X})\) 的分次（将习题 14 应用于 \(\sigma(x)\) 和 \(\sigma(y)\) 的最高次数非零齐次分量）；从而推出 \(\operatorname{Aut}(\mathrm{L}(\mathrm{X}))\) 同构于 \(\mathbf{GL}(2, \mathrm{K})\)。将这些结果推广到环 \(\mathrm{K}\) 没有幂零元 \(\neq 0\) 的情形。当 \(\mathrm{K}\) 含有平方为零的元素 \(\varepsilon \neq 0\) 时，证明 \(x \mapsto x, y \mapsto y + \varepsilon[x, y]\) 可以延拓为 \(\mathrm{L}(\mathrm{X})\) 的一个不保持 \(\mathrm{L}(\mathrm{X})\) 的分次的自同构。

¶ 16. 设 K 是一个 Q-代数。若 g 是一个李代数，用 U(g) 表示其包络代数，用 σ 表示 g 到 U(g) 的典范映射，并用 S(g) 表示 K-模 g 的对称代数。存在唯一的线性映射

\[
\eta (\mathfrak {g}): \mathrm{S} (\mathfrak {g}) \rightarrow \mathrm{U} (\mathfrak {g})
\]

满足 \(\eta (\mathfrak{g})(x^n) = \sigma (x)^n\) 对所有 \(x\in \mathfrak{g}\) 和所有 \(n\in \mathbf{N}\) 均成立。于是

\[
\eta (\mathfrak {g}) (x _ {1} \dots x _ {n}) = \frac {1}{n !} \sum_ {s \in \mathfrak {S} _ {n}} \sigma (x _ {s (1)} \dots x _ {s (n)}), \quad x _ {i} \in \mathfrak {g}.
\]

我们拟证明 \(\eta(g)\) 是双射。

\begin{enumerate}
\def\labelenumi{(\alph{enumi})}
\item
  证明 \(\eta(g)\) 是满射；当 K-模 \(g\) 是自由模时，证明它是双射（使用庞加莱–伯克霍夫–维特定理）。
\item
  令 \(\mathfrak{h}\) 是 \(\mathfrak{g}\) 的一个理想。令 \(\mathfrak{s}_{\mathfrak{h}}\)（相应地，\(\mathfrak{u}_{\mathfrak{h}}\)）表示 \(S(\mathfrak{g})\)（相应地，\(U(\mathfrak{g})\) 的双边理想）中由 \(\mathfrak{h}\)（相应地，\(\sigma(\mathfrak{h})\)）生成的理想。下图
\end{enumerate}

\[
\begin{array}{c} 0 \to s _ {\mathfrak {h}} \to S (g) \to S (g / \mathfrak {h}) \to 0 \\ \downarrow^ {\eta (g)} \quad \downarrow^ {\eta (g / \mathfrak {h})} \\ 0 \to u _ {\mathfrak {h}} \to U (g) \to U (g / \mathfrak {h}) \to 0 \end{array}
\]

图交换，且其各行正合。由此推出 \(\tilde{u}_{\mathfrak{h}}\)，即 \(s_{\mathfrak{h}}\) 在 \(\eta(g)\) 下的像，包含于 \(u_{\mathfrak{h}}\) 中；并且 \(\tilde{u}_{\mathfrak{h}} = u_{\mathfrak{h}}\)，只要 \(\eta(g)\) 和 \(\eta(g/h)\) 是单射（特别地，当模 \(g\) 和 \(g/h\) 是自由模时）。

\begin{enumerate}
\def\labelenumi{(\alph{enumi})}
\setcounter{enumi}{2}
\item
  取 \(\mathfrak{g}\) 为具有基本族 \((\mathbf{X}_1, \ldots, \mathbf{X}_n, \mathrm{H})\) 的自由李代数，其中 \(n \geqslant 0\)，并令 \(\mathfrak{h}\) 为由 H 生成的 \(\mathfrak{g}\) 的理想。证明 \(\mathfrak{g}\) 和 \(\mathfrak{g} / \mathfrak{h}\) 是自由模。由此推出在该情形下 \(\tilde{\mathfrak{u}}_{\mathfrak{h}} = \mathfrak{u}_{\mathfrak{h}}\)。特别地，存在 \(x \in \mathfrak{s}_{\mathfrak{h}}\)，使得 \((\eta(\mathfrak{g}))(x) = \sigma(\mathbf{X}_1) \ldots \sigma(\mathbf{X}_n) \sigma(\mathrm{H})\)。
\item
  回到一般情形。证明 \(\mathfrak{u}_{\mathfrak{h}}\) 作为 \(\mathrm{K}\)-模由形如 \(\sigma(x_1) \ldots \sigma(x_n)\sigma(h)\) 的元素生成，其中 \(n \geqslant 0\)、\(x_i \in \mathfrak{g}\) 且 \(h \in \mathfrak{h}\)（注意 \(\mathfrak{u}_{\mathfrak{h}}\) 与由 \(\mathfrak{h}\) 生成的左理想重合）。利用 (c) 证明这种元素属于 \(\tilde{\mathfrak{u}}_{\mathfrak{h}}\)（使用从自由李代数到 \(\mathfrak{g}\) 的适当同态）。推出 \(\tilde{\mathfrak{u}}_{\mathfrak{h}} = \mathfrak{u}_{\mathfrak{h}}\)。
\item
  证明若 \(\eta(g)\) 是双射，则 \(\eta(g/h)\) 也是双射。最后推出，对每个李代数 \(g, \eta(g)\) 都是双射（将 \(g\) 写成一个自由李代数的商），并且典范同态
\end{enumerate}

\[
\omega : \mathsf {S} (\mathfrak {g}) \to \operatorname{gr} \mathrm{U} (\mathfrak {g}) \quad (\text { cf.   Chapter   I,   } \S 2, \text {   no.   } 6)
\]

是同构（Q-代数上李代数的“庞加莱–伯克霍夫–维特定理”）。

\exercisesfor{3}

字母 X 表示一个集合。

\begin{enumerate}
\def\labelenumi{\arabic{enumi}.}
\tightlist
\item
  \begin{enumerate}
  \def\labelenumii{(\alph{enumii})}
  \tightlist
  \item
    证明每个元素 \(u \in \mathrm{A}^{+}(\mathrm{X})\) 都能唯一地写成 \(u = \sum_{x \in \mathrm{X}} u_{x} x\) 的形式，其中 \(u_{x} \in \mathrm{A}(\mathrm{X})\)。
  \end{enumerate}
\end{enumerate}

\begin{enumerate}
\def\labelenumi{(\alph{enumi})}
\setcounter{enumi}{1}
\tightlist
\item
  证明 \(\{0\}\) 是 \(\mathrm{A}^{+}(\mathrm{X})\) 中唯一在所有映射 \(u \mapsto u_{x} (x \in \mathrm{X})\) 下稳定的子模。（若 a 是这样的子模且 \(a \neq \{0\}\)，考虑 a 中次数最小的非零元素。）
\end{enumerate}

\begin{enumerate}
\def\labelenumi{\arabic{enumi}.}
\setcounter{enumi}{1}
\tightlist
\item
  设 g 是一个自由模李代数；通过 \(\sigma: g \to Ug\)，将其与它在包络代数 \(Ug\) 中的像相识别。令 \(U^{+}g\) 为典范同态 \(Ug \to K\) 的核。令 i 是从 X 到 g 的映射，使得 \(i(X)\) 作为李代数生成 g。
\end{enumerate}

\begin{enumerate}
\def\labelenumi{(\alph{enumi})}
\item
  证明 \(\mathbf{U}^{+}\mathfrak{g}\) 作为左 Ug-模由 \(i(\mathbf{X})\) 生成。
\item
  证明以下性质等价：
\item
  g 是以 i: X → g 为基本族的自由李代数。
\end{enumerate}

\begin{enumerate}
\def\labelenumi{(\roman{enumi})}
\setcounter{enumi}{1}
\tightlist
\item
  族 \(i\) 是左 Ug-模 \(\mathbf{U}^{+}\mathfrak{g}\) 的一组基。
\end{enumerate}

（蕴含关系 (i) \(\Rightarrow\) (ii) 利用同构 \(\mathrm{Ug} \to \mathrm{A}(\mathrm{X})\)，由习题 1 (a) 得出。为证明 (ii) \(\Rightarrow\) (i)，将习题 1 (b) 应用于同态 \(\mathrm{A}(\mathrm{X}) \to \mathrm{Ug}\) 的核，该同态由 \(i\) 定义。）

\begin{enumerate}
\def\labelenumi{\arabic{enumi}.}
\setcounter{enumi}{2}
\tightlist
\item
  设 \(x \in \mathbf{X}\)。证明 \(x\) 在 \(\mathrm{A}(\mathbf{X})\) 中的中心化子是由 \(x\) 生成的子代数。由此推出，\(\mathrm{L}(\mathbf{X})\) 中与 \(x\) 对易的唯一元素是 \(x\) 的倍数。特别地，\(\mathrm{L}(\mathbf{X})\) 的中心在 \(\operatorname{Card}(\mathbf{X}) \geqslant 2\) 时为 (0)，而 \(\mathrm{L}(\mathbf{X}) / \left( \sum_{n > p} \mathrm{L}^n(\mathbf{X}) \right)\) 的中心化为 \(\mathrm{L}^p(\mathbf{X})\) 的典范像。
\end{enumerate}


¶ 4. 设 K 是特征为 p \textgreater{} 0 的域。

\begin{enumerate}
\def\labelenumi{(\alph{enumi})}
\item
  设 \(\mathbf{H}\) 是 \(\mathbf{L}(\mathbf{X})\)（它是 \(\mathbf{A}(\mathbf{X})\) 的李子代数）的一组基。证明（利用同构 \(\mathbf{U}(\mathbf{L}(\mathbf{X})) \to \mathbf{A}(\mathbf{X})\) 和庞加莱–伯克霍夫–维特定理），元素 \(h^{p^n}\)（\(h \in \mathbf{H}\)、\(n \in \mathbf{N}\)）在 \(\mathbf{K}\) 上线性无关。
\item
  若 \(m\) 是整数 \(\geqslant 0\)，令 \(L_m\) 表示 \(A(X)\) 中以 \(h^{p^n}\)（其中 \(h \in H, n \leqslant m\)）为基的子模。证明 \(L_m\) 是 \(A(X)\) 的李子代数，并且若 \(a \in L_m - 1\)，则 \(a^p \in L_m\)（对 \(m\) 归纳，并使用 Jacobson 公式，参见~第一章，§ 1，习题 19）。推出 \(L_m\) 与 \(H\) 的选择无关，并且它是 \(A(X)\) 中由 \(x^{p^n}\)（其中 \(x \in X, n \leqslant m\)）生成的李子代数。
\item
  令 \(\mathbf{L}(\mathbf{X}, p)\) 为各 \(\mathbf{L}_m\)（\(m \geqslant 0\)）的并。证明 \(\mathbf{L}(\mathbf{X}, p)\) 是最小李 \(p\)-子代数，包含于 \(\mathbf{A}(\mathbf{X})\) 且包含 \(\mathbf{X}\)（参见~第一章 §1 的习题 20）；它以 \(h^{p^n}\)（其中 \(h \in \mathbf{H}\)、\(n \in \mathbf{N}\)）为基。
\item
  令 \(\tilde{\mathbf{U}}(\mathbf{L}(\mathbf{X},p))\) 为 \(\mathbf{L}(\mathbf{X},p)\) 的受限包络代数，参见~第一章 §2 的习题 6。证明 \(\mathbf{L}(\mathbf{X},p)\) 到 \(\mathbf{A}(\mathbf{X})\) 的嵌入可以延拓为 \(\tilde{\mathbf{U}}(\mathbf{L}(\mathbf{X},p))\) 到 \(\mathbf{A}(\mathbf{X})\) 的同构。（使用庞加莱–伯克霍夫–维特定理和上面引用的习题。）推出 \(\mathbf{L}(\mathbf{X},p)\) 是 \(\mathbf{A}(\mathbf{X})\) 的本原元素集合，参见~§1 的习题 12。
\item
  设 \(f\) 是从 \(\mathbf{X}\) 到李 \(p\)-代数 \(\mathfrak{g}\) 的映射。证明 \(f\) 可以唯一延拓为 \(p\)-同态 \(\mathbf{F} \colon \mathbf{L}(\mathbf{X}, p) \to \mathfrak{g}\)。（先将 \(f\) 延拓为从 \(\mathrm{A}(\mathbf{X})\) 到 \(\tilde{\mathrm{U}}\mathfrak{g}\) 的代数同态。）
\end{enumerate}

（李 \(p\)-代数 \(\mathbf{L}(\mathbf{X}, p)\) 称为自由李 \(p\)-代数（集合 \(\mathbf{X}\) 上）。）

\exercisesfor{4}

在下列习题中，字母 G 表示一个群。

\begin{enumerate}
\def\labelenumi{\arabic{enumi}.}
\item
  设 \((\mathbf{G}_n)\) 是 \(G\) 上的一个整中心滤过。证明李代数 \(\mathrm{gr}(G)\) 由 \(\mathrm{gr}_1(G)\) 生成，当且仅当 \(\mathbf{G}_n = \mathbf{G}_{n+1}.\mathbf{C}^n\mathbf{G}\) 对所有 \(n \geqslant 1\) 成立。在此情形，证明 \(\mathbf{G}_n = \mathbf{G}_m.\mathbf{C}^n\mathbf{G}\) 对 \(m > n\) 成立，并推出 \((\mathbf{G}_n) = (\mathbf{C}^n\mathbf{G})\)，若存在整数 \(m\) 使 \(\mathbf{G}_m = \{\varepsilon\}\)。
\item
  设 \((\mathbf{G}_{\alpha})\) 是群 \(\mathbf{G}\) 上的实滤过，\(v\) 是相应的阶函数。令 \(\mathbf{H}\) 是 \(\mathbf{G}\) 的一个子群。
\end{enumerate}

\begin{enumerate}
\def\labelenumi{(\alph{enumi})}
\item
  令 \(H_{\alpha} = H \cap G_{\alpha}\)。证明 \((\mathrm{H}_{\alpha})\) 是 H 上的实滤过，且相应的阶函数是 v 在 H 上的限制。\(H_{\alpha}^{+} = H \cap G_{\alpha}^{+}\)，并且 \(\mathrm{gr}(\mathrm{H})\) 可与 \(\mathrm{gr}(\mathrm{G})\) 的一个分次子群相识别。
\item
  假设 H 是正规子群，且 \(v(G) \cap \mathbf{R}\) 是 R 的离散子集。记 \((G/H)_{\alpha} = (G_{\alpha}H)/H\)。证明 \(((G/H)_{\alpha})\) 是 G/H 上的实滤过，且相应的阶函数 \(v_{G/H}\) 由下式给出
\end{enumerate}

\[
v _ {\mathrm{G/H}} (x) = \operatorname * {S u p} _ {y \in x} v (y).
\]

证明对所有 \(\alpha \in \mathbf{R}\)，存在正合列

\[
0 \rightarrow \mathrm{gr} _ {\alpha} (\mathrm{H}) \rightarrow \mathrm{gr} _ {\alpha} (\mathrm{G}) \rightarrow \mathrm{gr} _ {\alpha} (\mathrm{G} / \mathrm{H}) \rightarrow 0.
\]

\begin{enumerate}
\def\labelenumi{(\alph{enumi})}
\setcounter{enumi}{2}
\tightlist
\item
  在 \((b)\) 的假设下，假设 \((G_{\alpha})\) 是中心滤过。证明由 \((G_{\alpha})\) 在 H 和 G/H 上诱导的滤过也都是中心滤过，并且李代数 \(\mathrm{gr}(\mathrm{G}/\mathrm{H})\) 可与 \(\mathrm{gr}(\mathrm{G})\) 对理想 \(\mathrm{gr}(\mathrm{H})\) 取商所得的代数相识别。
\end{enumerate}

¶ 3. 设 \((\mathrm{H}_{\alpha})\) 是群 H 上的实滤过，\(v_{\mathrm{H}}\) 是相应的阶函数。假设 G 在 H 上左作用，并且对所有 \(g \in G\)，映射 \(h \mapsto g(h)\) 都是滤过群 H 的自同构。

\begin{enumerate}
\def\labelenumi{(\alph{enumi})}
\tightlist
\item
  若 \(\alpha \in \mathbf{R}\)，用 \(G_{\alpha}\) 表示满足下列条件的 \(g \in G\) 的集合：
\end{enumerate}

\[
v _ {\mathrm{H}} (h ^ {- 1} g (h)) \geqslant v _ {\mathrm{H}} (h) + \alpha \quad \text { for   all } h \in \mathrm{H}.
\]

证明 \((\mathbf{G}_{\alpha})\) 是 \(\mathbf{G}\) 上的实滤过，并且 \(\mathbf{G}_{\alpha} = \mathbf{G}\) 当 \(\alpha \leqslant 0\) 时。令 \(v\) 表示相应的阶函数。

\begin{enumerate}
\def\labelenumi{(\alph{enumi})}
\setcounter{enumi}{1}
\tightlist
\item
  假设滤过 \((\mathrm{H}_{\alpha})\) 是中心滤过，且 \(\mathrm{G}_0^+ = \mathrm{G}\)。证明 \((\mathrm{G}_{\alpha})\) 是 G 上的中心滤过。若 \(\xi \in \mathrm{gr}_{\alpha}(\mathrm{G})\)、\(\eta \in \mathrm{gr}_{\beta}(\mathrm{H})\)，令 \(g\)（相应地，\(h\)）为 \(\xi\)（相应地，\(\eta\)）在 \(\mathrm{G}_{\alpha}\)（相应地，\(\mathrm{H}_{\beta}\)）中的代表元；证明 \(h^{-1}g(h)\) 在 \(\mathrm{gr}_{\alpha + \beta}(\mathrm{H})\) 中的像与 \(g\)、\(h\) 的选择无关；若将其记为 \(\mathrm{D}_{\xi}(\eta)\)，证明 \(\mathrm{D}_{\xi}\) 可延拓为 \(\mathrm{gr}(\mathrm{H})\) 的次数为 \(\alpha\) 的一个导子，并且 \(\xi \mapsto \mathrm{D}_{\xi}\) 是从李代数 \(\mathrm{gr}(\mathrm{G})\) 到 \(\mathrm{gr}(\mathrm{H})\) 的导子李代数的一个同态。若 \(v_{\mathrm{H}}(\mathrm{H}) \cap \mathbf{R}\) 包含于 \(\Gamma\)（\(\mathbf{R}\) 的离散子群），则 \(v(\mathrm{G}) \cap \mathbf{R}\) 也包含于该子群，并且对所有 \(\alpha \in \mathbf{R}\)，上面定义的映射 \(\xi \mapsto \mathrm{D}_{\xi}\) 都是单射。
\end{enumerate}

\begin{enumerate}
\def\labelenumi{\arabic{enumi}.}
\setcounter{enumi}{3}
\item
  设 H 是类为 c 的幂零群，G 是在 \(\mathrm{H}/(\mathrm{H},\mathrm{H})\) 上平凡作用的 H 的自同构群。证明 G 是类 \(\leqslant c - 1\) 的幂零群。（将习题 3 应用于 H 的下中心列，并注意 G 在 \(\mathrm{gr}(\mathrm{H})\) 上平凡作用。）证明若 H 是有限 p-群（p 为素数），则 G 也是如此（方法相同）。
\item
  设 \(\mathbf{K}\) 是一个交换域，\(\mathbf{L}\) 是 \(\mathbf{K}\) 的有限 Galois 扩张，Galois 群为 \(\mathbf{G}\)。令 \(v_{\mathbf{L}}\) 是 \(\mathbf{L}\) 的一个取值于 \(\mathbf{Z}\) 的赋值（交换代数，第六章，§3，第 2 号），且在 \(\mathbf{G}\) 下不变。若 \(g \in \mathbf{G}\)，记
\end{enumerate}

\[
v (g) = \sup _ {x \in L ^ {\bullet}} v _ {L} \left(\frac {g (x) - x}{x}\right).
\]

证明 v 是 G 上一个分离的整滤过 \((G_{n})\) 的阶函数，满足 \(G_{0} = G\)，且此滤过在 \(G_{1}\) 上的限制是中心滤过（应用习题 3，取 H 为 \(v_{L}\) 的理想）。证明 \(G_{1}\) 是 p-群（相应地，化为单位元），当 L 的剩余域的特征为 p \textgreater{} 0（相应地，为零）时；当域 L 和 K 具有相同的剩余域时，\(G_{1}\) 是交换代数第六章 §8 习题 11 (b) 中定义的同态 \(\phi\) 的核。

¶ 6. 令 K{[}G{]} 为 G 在 K 上的代数，I 为典范同态 K{[}G{]} → K 的核（“增广理想”）。于是


\(K[G] = K \oplus I\)，且 \(I\) 以 \(g - 1\) 为基，其中 \(g \in G - \{e\}\)。

\begin{enumerate}
\def\labelenumi{(\alph{enumi})}
\item
  若 \(n\) 是整数 \(\geqslant 0\)，令 \(I^n\) 表示 \(K[G]\) 的第 \(n\) 次幂 \(I\) 所成的理想。令 \(G_n\) 为由满足 \(g \in G\) 且 \(g - 1 \in I^n\) 的元素组成的集合。证明 \((G_n)\) 是 \(G\) 上的整中心滤过。特别地，\(G_n \supset C^n G\) 对所有 \(n\) 成立。
\item
  证明当 \(\mathbf{K} = \mathbf{Z}\) 时，映射 \(g \mapsto g - 1\) 在取商后定义了从 \(\mathrm{G} / (\mathrm{G}, \mathrm{G})\) 到 \(\mathrm{I} / \mathrm{I}^2\) 的同构。推出 \(\mathrm{G}_2 = \mathrm{C}^2\mathrm{G}.\dagger\)
\item
  设 K 是特征为零的域且 G 有限。证明 \(I^{n} = I\) 对所有 \(n \geqslant 1\) 成立。
\item
  设 K 是特征为 p \textgreater{} 0 的域且 G 是 p-群。证明当 n 充分大时 \(I^{n}\{0\}\)。（首先利用代数第一章 §6 第 5 号命题 11 证明每个单 K{[}G{]}-模都同构于 K；推出 I 是 K{[}G{]} 的根基，因而是幂零的，因为 K{[}G{]} 在 K 上有限秩。）
\item
  设 K = Z，且 G 具有如下性质：对每个 \(g \in G\)，若 \(g \neq 1\)，则存在素数 p、一个 p-群 P 以及同态 \(f: G \to P\)，使得 \(f(g) \neq e\)。证明在此情形下 \(\bigcap_{n} I^{n} = \{0\}\)。（立即将其归结为 G = P 的情形。将 (d) 应用于域 \(F_{p}\)，可知存在 m 使 \(I^{m} \subset p\)。\(Z[G]\)，而 I 是 Z{[}G{]} 中的直因子，所以这意味着 \(I^{m} \subset pI\)，从而 \(\bigcap_{n} I^{mn} \subset \bigcap_{n} p^{n}I\)；后者为 0，因为 I 是有限生成的 Abel 群。）
\end{enumerate}

\begin{enumerate}
\def\labelenumi{\arabic{enumi}.}
\setcounter{enumi}{6}
\item
  给 G 赋予滤过 \((\mathrm{C}^{\mathrm{n}}\mathrm{G})\)，并假设 \(\mathrm{gr}_{1}(\mathrm{G}) = \mathrm{G}/(\mathrm{G}, \mathrm{G})\) 是循环群。证明 \(\mathrm{gr}_{n}(\mathrm{G}) = \{0\}\) 对 \(n \geqslant 2\) 成立（使用命题 5），并推出 \(\mathrm{C}^{\mathrm{n}}\mathrm{G} = (\mathrm{G}, \mathrm{G})\) 对 \(n \geqslant 2\) 成立。
\item
  令 \(G = \mathbf{S}\mathbf{L}_{2}(\mathbf{Z})\)，并令 \(x = \begin{pmatrix} 1 & 1 \\ 0 & 1 \end{pmatrix}, y = \begin{pmatrix} 1 & 0 \\ 1 & 1 \end{pmatrix}, w = \begin{pmatrix} 0 & 1 \\ -1 & 0 \end{pmatrix}.\)
\end{enumerate}

\begin{enumerate}
\def\labelenumi{(\alph{enumi})}
\item
  验证公式 \(w^4 = 1\)、\(w = xy^{-1}x\)、\(wxw^{-1} = y^{-1}\)。
\item
  若 \(g = \begin{pmatrix} a & b \\ c & d \end{pmatrix}\) 是 \(G\) 的元素，记 \(l(g) = |a| + |b|\)。证明 \(l(g) = 1\) 当且仅当 \(g\) 具有 \(y^n w^a\) 的形式，其中 \(n \in \mathbf{Z}\) 且 \(0 \leqslant a \leqslant 3\)。若 \(l(g) \geqslant 2\)，证明存在某个幂 \(h\)（它是 \(x\) 或 \(y\) 的幂），使 \(l(gh) < l(g)\)。推出 \(G\) 由 \(\{x, y\}\) 生成。
\item
  利用 (a) 和 (b)，证明 G/(G, G) 由 x 的像 \(\xi\) 生成，且 \(\xi^{12} = e.\ddagger\) 推出 \(C^{n}G = (G, G)\) 对 \(n \geqslant 2\) 成立。（应用习题 7。）
\end{enumerate}

\begin{enumerate}
\def\labelenumi{\arabic{enumi}.}
\setcounter{enumi}{8}
\tightlist
\item
  给 G 赋予滤过 \((C^{n}G)\)。
\end{enumerate}

\begin{enumerate}
\def\labelenumi{(\alph{enumi})}
\item
  证明 \(\operatorname{gr}(\mathbf{G})\) 中由 \(\operatorname{gr}_2(\mathbf{G})\) 生成的理想是 \(\sum_{q\geqslant 2}\operatorname{gr}_q(\mathbf{G})\)。
\item
  令 \((x_{i})_{i\in I}\) 是 \(G\) 的生成族，并令 \(m\geqslant 1\)。假设对所有 \((i,j)\in \mathrm{I}^2\)，都有 \((x_i,x_j)^m\in \mathrm{C}^3\mathrm{G}\)。证明对所有 \(q\geqslant 2\) 及所有 \(u\in \mathrm{C}^q\mathrm{G}\)，都有 \(u^{m}\in \mathrm{C}^{q + 1}\mathrm{G}\)（使用 \((a)\)）。
\end{enumerate}

\begin{enumerate}
\def\labelenumi{\arabic{enumi}.}
\setcounter{enumi}{9}
\tightlist
\item
  设 \(x, y\) 属于 \(G\)，并令 \(r, s\) 是两个整数 \(\geqslant 1\)。假设 \((x^r, y^s) = e\)。（a）证明对所有 \(n \geqslant 1\)，都有 \((x, y)^{(rs)n} \in C^{n+2}G\)。（可以假设 \(G\) 由 \(\{x, y\}\) 生成；然后应用习题 9 (b)，注意 \((x^r, y^s) \equiv (x, y)^{rs} \mod C^3G\)。）
\end{enumerate}

\begin{enumerate}
\def\labelenumi{(\alph{enumi})}
\setcounter{enumi}{1}
\tightlist
\item
  假设 \(x^r = y^s = e\)；令 \(t\) 表示 \(r\) 和 \(s\) 的最大公约数。证明 \((x, y)^t \in \mathbf{C}^3\mathbf{G}\)，并推出 \((x, y)^{t^n} \in \mathbf{C}^{n+2}\mathbf{G}\) 对所有 \(n \geqslant 1\) 成立（方法相同）。
\end{enumerate}

\begin{enumerate}
\def\labelenumi{\arabic{enumi}.}
\setcounter{enumi}{10}
\tightlist
\item
  设 H 是 G 的一个子群，m 是满足 \(\geqslant 1\) 的整数。假设 G 由一个族 \((x_{i})_{i\in I}\) 生成，且对所有 i 都有 \(x_{i}^{m}\in H\)。
\end{enumerate}

\begin{enumerate}
\def\labelenumi{(\alph{enumi})}
\tightlist
\item
  令 H 带有由 G 上的滤过 (C \(^{n}\) G) 诱导的滤过，并令 gr(H) 与 gr(G) 的一个分次李子代数相识别，参见~习题 2。证明对所有 \(n \geqslant 0\)，
\end{enumerate}

\[
m ^ {n}. \mathrm{gr} _ {n} (\mathrm{G}) \subset \mathrm{gr} _ {n} (\mathrm{H}).
\]

推出对所有 \(z \in \mathbf{H}.\mathbf{C}^n\mathbf{G}\)，都有 \(z^m \in \mathbf{H}.\mathbf{C}^{n+1}\mathbf{G}\)。

\begin{enumerate}
\def\labelenumi{(\alph{enumi})}
\setcounter{enumi}{1}
\tightlist
\item
  若 G 是幂零群，证明存在整数 \(N \geqslant 0\)（只取决于 G 的幂零类），使得 \(z^{mN} \in H\) 对所有 \(z \in G\) 成立。
\end{enumerate}

\begin{enumerate}
\def\labelenumi{\arabic{enumi}.}
\setcounter{enumi}{11}
\item
  假设 G 是幂零群。令 H 是 G 的一个子群，m 是满足 \(\geqslant 1\) 的整数，且令 x、y 是 G 中满足 \(x^{m} \in H\) 和 \(y^{m} \in H\) 的元素。证明存在整数 \(N \geqslant 0\)，使得 \((xy)^{nN} \in H\)。（将习题 11 应用于由 \(\{x, y\}\) 生成的群及其与 H 的交。）
\item
  \begin{enumerate}
  \def\labelenumii{(\alph{enumii})}
  \tightlist
  \item
    令 \(\mathbf{F}\) 是具有两个元素的基本族 \(\{\overline{x},\overline{y}\}\) 的自由群，令 \(c\) 是整数 \(\geqslant 2\)，并令 \(m\) 是整数 \(\geqslant 1\)。记 \(\mathrm{F}^c = \mathrm{F} / C^c\mathrm{F}\)，并令 \(x,y\) 为 \(\overline{x},\overline{y}\) 在 \(\mathrm{F}^c\) 中的像。令 \(\mathrm{F}_m^c\) 为 \(\mathrm{F}^c\) 中由 \(\{x^{m},y\}\) 生成的子群。证明存在整数 \(\mathrm{N}\geqslant 0\)，使 \(z^{mN}\in \mathrm{F}_{m}^{c}\) 对所有 \(z\in \mathrm{F}^c\) 成立（使用习题 11）。
  \end{enumerate}
\end{enumerate}

\begin{enumerate}
\def\labelenumi{(\alph{enumi})}
\setcounter{enumi}{1}
\item
  令 \(\mathrm{I}^c\)（相应地，\(\mathbf{I}_m^c\)）为 \(\mathrm{F}^c\)（相应地，\(\mathbf{F}_m^c\)）的由 \(y\) 生成的正规子群。证明若 \(N\) 如上选取且 \(z \in \mathrm{I}^c\)，则 \(z^{n^k} \in \mathbf{I}_m^c\)（注意 \(\mathrm{F}^c / \mathrm{I}^c\) 是无限循环群，其生成元为 \(x\) 的像，并推出 \(\mathbf{F}_m^c / \mathbf{I}_m^c \to \mathbf{F}^c / \mathrm{I}^c\) 是单射，从而 \(\mathbf{I}_m^c = \mathbf{F}_m^c \cap \mathbf{I}^c\)）。特别地，\(xy^m x^{-1} \in \mathbf{I}_m^c\)。
\item
  假设 G 是幂零群。令 H 是 G 的一个子群，L 是 H 的一个正规子群。令 \(g \in G\)，并令 \(m \geqslant 1\) 使得 \(g^{m} \in H\)。证明当 N 充分大时，\(gl^{m^{N}} g^{-1} \in L\) 对所有 \(l \in L\) 成立。（若 G 的类 \textless c，取 (a) 中的 N，并使用同态 \(f: F^{c} \to G\)，满足 \(f(x) = g, f(y) = l\)；注意 \(f(F_{m}^{c}) \subset H, f(I_{m}^{c}) \subset L\)，然后应用上面的 (b)。）
\end{enumerate}

¶ 14. 设 P 是素数集合。若整数 n 满足 n ≠0 且其所有素因子都属于 P，则称 n 为 P-整数。若元素 \(x \in G\) 满足存在 P-整数 n 使 \(x^{n} = \epsilon\)，则称其为 P-挠元素；若 G 的每个元素都是 P-挠元素（相应地，除 e 外 G 中没有 P-挠元素），则称 G 为 P-挠群（相应地，P-无挠群）。若对所有 \(x \in G\) 和每个 P-整数 n，都存在 \(y \in G\) 使 \(x = y^{n}\)，则称 G 为 P-可除群。

假设 G 是幂零群。

\begin{enumerate}
\def\labelenumi{(\alph{enumi})}
\item
  令 H 是 G 的一个子群，令 \(H_{P}\) 为满足对 \(x \in G\) 存在 P-整数 n 使 \(x^{n} \in H\) 的元素所成的集合。证明 \(H_{P}\) 是 G 的子群（使用习题 12），且 \((\mathrm{H}_{\mathrm{P}})_{\mathrm{P}} = \mathrm{H}_{\mathrm{P}}\)。群 \(H_{P}\) 称为 H 在 G 中的 P-饱和。若 G 是 P-可除群，则 \(H_{P}\) 也是 P-可除群。
\item
  令 \(\mathbf{L}\) 是 \(\mathbf{H}\) 的正规子群。证明 \(\mathbf{L}_{\mathbb{P}}\) 是 \(\mathbf{H}_{\mathbb{P}}\) 的正规子群。（使用习题 13 (c) 证明：若 \(g \in \mathbf{H}_{\mathbb{P}}\) 且 \(l \in \mathbf{L}\)，则 \(glg^{-1} \in \mathbf{L}_{\mathbb{P}}\)。）
\end{enumerate}

特别地，若 L 是 G 的正规子群，则 \(L_{P}\) 也是正规子群，且 \(G/L_{P}\) 是 P-无挠的。推出 G 的 P-挠元素集合是使 G/N 为 P-无挠群的最小正规子群 N。

\begin{enumerate}
\def\labelenumi{(\alph{enumi})}
\setcounter{enumi}{2}
\item
  假设 G 是 P-无挠的。令 \(n\) 是 P-整数，且 \(x, y \in G\) 满足 \(x^n = y^n\)。证明 \(x = y\)。（对 \(r = s = n\) 应用习题 10 (a)。推出存在 P-整数 N 使 \((x, y)^{\mathbb{N}} = e\)，于是 \((x, y) = e\)，因为 G 是 P-无挠的；然后 \((x^{-1}y)^n = e\)，最终得到 \(x = y\)。）
\item
  令 H 是 G 的一个子群，L 是一个幂零群，且 \(f: H \to L\) 是同态。令 \(\Gamma\) 为 f 在 \(H \times L\) 中的图，令 \(\Gamma_{P}\) 为它在 \(G \times L\) 中的 P-饱和。证明 \(\Gamma_{P}\) 包含于 \(H_{P} \times L\)，并且当 L 是 P-可除群时 \(pr_{1}: \Gamma_{P} \to H_{P}\) 是满射，当 L 是 P-无挠群时是单射。
\end{enumerate}

假设 L 是 P-可除且 P-无挠的。证明 f 可唯一延拓为同态 \(f_{P}: H_{P} \to L\)，且 \(f_{P}\) 的图就是 \(\Gamma_{P}\)。

¶ 15. 沿用上题的记号，假设 G 是幂零群。令 \(i: \mathbf{G} \to \overline{\mathbf{G}}\) 是从 G 到幂零群 \(\overline{\mathbf{G}}\) 的同态。若满足下列条件，则称 \((i, \overline{\mathbf{G}})\) 为 G 的 P-包络：

\begin{enumerate}
\def\labelenumi{(\roman{enumi})}
\item
  \(\overline{\mathbf{G}}\) 是 P-可除且 P-无挠的。
\item
  \(i\) 的核是 G 的 P-挠元素集合。
\item
  \(i(\mathbf{G})\) 在 \(\overline{\mathbf{G}}\) 中的 P-饱和等于 \(\overline{\mathbf{G}}\)。
\end{enumerate}

\begin{enumerate}
\def\labelenumi{(\alph{enumi})}
\item
  令 \((i, \overline{\mathbf{G}})\) 是 G 的一个 P-包络，L 是 P-可除、P-无挠的幂零群。证明对每个同态 \(f\colon \mathbf{G} \to \mathbf{L}\)，存在唯一的同态 \(\bar{f}\colon \overline{\mathbf{G}} \to \mathbf{L}\)，满足 \(\bar{f} \circ i = f\)（将其归结为 G 是 P-无挠的情形，并使用习题 14 (d)）。推出若 G 有 P-包络，\(\dagger\) 则它在同构意义下唯一。
\item
  令 \((i, \overline{\mathbf{G}})\) 是 G 的一个 P-包络，令 H 是 G 的一个子群，令 \(\overline{\mathbf{H}}\) 是 \(i(\mathbf{H})\) 在 \(\overline{\mathbf{G}}\) 中的 P-包络，并令 \(i_{\mathbf{H}}: \mathbf{H} \to \overline{\mathbf{H}}\) 是由 \(i\) 诱导的同态。证明 \((i_{\mathbf{H}}, \overline{\mathbf{H}})\) 是 H 的一个 P-包络。
\end{enumerate}

假设 H 是 G 的正规子群。则 \(\bar{H}\) 是 \(\bar{G}\) 的正规子群，参见~习题 14 (b)；若 \(i_{G/H}: G/H \to \overline{G}/\overline{H}\) 是由 i 诱导的同态，证明 \((i_{G/H}, \overline{G}/\overline{H})\) 是 G/H 的一个 P-包络。

\begin{enumerate}
\def\labelenumi{\arabic{enumi}.}
\setcounter{enumi}{15}
\tightlist
\item
  沿用前两题的记号。
\end{enumerate}

\begin{enumerate}
\def\labelenumi{(\alph{enumi})}
\item
  令 \(\mathbf{S}_{\mathbf{P}}\) 为 P-整数的集合，令 \(\mathbf{\dot{Z}}_{\mathbf{P}} = \mathbf{S}_{\mathbf{P}}^{-1}\mathbf{Z}\) 为 \(\mathbf{Z}\) 关于 \(\mathbf{S}_{\mathbf{P}}\) 定义的分式环（交换代数，第二章，§2，第 1 号）。假设 G 是幂零、P-无挠且 P-可除的，并令 \(t \in \mathbf{Z}_{\mathbf{P}}\)、\(g \in G\)。证明存在唯一的属于 G 的元素 \(h\)，使得 \(h^s = g^{st}\) 对所有 \(s \in \mathbf{S}_{\mathbf{P}}\) 且 \(st \in \mathbf{Z}\) 的 s 成立。元素 \(h\) 称为 \(t\) 次 \(g\) 的幂，记为 \(g^t\)。映射 \(t \mapsto g^t\) 是从 \(\mathbf{Z}_{\mathbf{P}}\) 到 G 的同态。若 \(t\) 在 \(\mathbf{Z}_{\mathbf{P}}\) 中可逆，则 \(g \mapsto g^t\) 是双射。（使用习题 14 (c)。）
\item
  令 \(\mathbf{A}\) 是交换群，令 \(i\) 为从 \(\mathbf{A}\) 到 \(\mathbf{A}_{\mathbb{P}} = \mathbf{Z}_{\mathbb{P}} \times \mathbf{A}\) 的典范映射。证明 \((i, \mathbf{A}_{\mathbb{P}})\) 是 \(\mathbf{A}\) 的一个 P-包络。
\item
  令 \(\mathbf{G}\)（相应地，\(\overline{\mathbf{G}}\)）为 \(n\) 阶的环 \(k\)（相应地，环 \(k_{\mathrm{P}} = k\otimes \mathbf{Z}_{\mathrm{P}}\)）上的严格下三角群，并令 \(i\) 为从 \(\mathbf{G}\) 到 \(\overline{\mathbf{G}}\) 的由 \(k\) 到 \(k_{\mathrm{P}}\) 的典范同态所定义的同态（代数，第二章，§10，第 4 号）。证明 \((i,\overline{\mathbf{G}})\) 是 \(\mathbf{G}\) 的一个 P-包络。
\item
  令 G 是有限幂零群，因而是其 Sylow \(p\)-群 \(G_p\) 的直积（代数，第一章，§6，第 7 号，定理 4）。令 \(\overline{G}\) 为 \(G_p\)（\(p \notin P\)）的乘积，并令 \(i\) 为从 G 到 \(\overline{G}\) 的典范投影。证明 \((i, \overline{G})\) 是 G 的一个 P-包络。
\end{enumerate}

¶ 17. 假设 G 是幂零群。令 P 是素数集合，令 \(\mathfrak{V}_{P}(G)\) 为 G 中有限指数且指数为 P-整数的正规子群的集合（参见~习题 14）。令 \(\mathcal{T}_{P}(G)\) 为由滤子基 \(\mathfrak{V}_{P}(G)\) 在 G 上定义的拓扑（一般拓扑，第三章，§1，第 2 号，例）。

\begin{enumerate}
\def\labelenumi{(\alph{enumi})}
\item
  令 N 是 G 的一个有限指数子群，且 (G: N) 是 P-整数，令 \(N'\) 为 N 的所有共轭子群的交。证明 \((G: N')\) 是 P-整数（利用群 \(G/N'\) 是其 Sylow 子群的乘积这一事实）；推出 N 关于 \(\mathcal{T}_{\mathbb{P}}(G)\) 是开的。
\item
  令 H 是 G 的有限指数子群。证明 \(\mathcal{T}_{\mathrm{P}}(\mathrm{H})\) 与 \(\mathcal{T}_{\mathrm{P}}(\mathrm{G})\) 在 H 上诱导的拓扑重合。（方法相同。）
\item
  令 H 是 G 的正规子群，且 G/H 同构于 Z，令 x 是 G/H 的一个生成元在 G 中的代表元。令 \(N \in \mathfrak{V}_{\mathrm{P}}(\mathrm{H})\)，并令 \(N' = \bigcap_{n \in Z} x^{n} N x^{-1}\)。证明 \(N' \in \mathfrak{V}_{\mathrm{P}}(\mathrm{H})\)，且 \(x N' x^{-1} = N'\)。若 \(N''\) 是由 \(N'\) 和 x 生成的子群，证明 \((G: N'') = (H: N')\)。推出 \(\mathcal{T}_{\mathrm{P}}(\mathrm{H})\) 是由 \(\mathcal{T}_{\mathrm{P}}(\mathrm{G})\) 诱导的。
\item
  假设 \(\mathbf{G}\) 是有限群。证明若 \(\mathbf{H}\) 是 \(\mathbf{G}\) 的子群，则存在子群序列
\end{enumerate}

\[
\mathrm{H} = \mathrm{H} _ {0} \subset \mathrm{H} _ {1} \subset \dots \subset \mathrm{H} _ {k} = \mathrm{G}
\]

满足 \(\mathbf{H}_i\) 是 \(\mathrm{H}_{i + 1}\) 的正规子群，对 \(0\leqslant i < k\) 成立，且 \(\mathrm{H}_{i + 1} / \mathrm{H}_i\) 是有限群或

同构于 \(\mathbf{Z}\)。利用 \((b)\) 和 \((c)\)，推出 \(\mathcal{T}_{\mathrm{P}}(\mathrm{H})\) 是由 \(\mathcal{T}_{\mathrm{P}}(\mathrm{G})\) 诱导的。

\begin{enumerate}
\def\labelenumi{(\alph{enumi})}
\setcounter{enumi}{4}
\tightlist
\item
  在 (d) 的假设下，令 P' 表示素数集合中 P 的补集。证明 H 关于 \(\mathcal{T}_{\mathrm{P}}(\mathrm{G})\) 的闭包与 H 在 G 中的 P'-饱和重合（习题 14）；特别地，G 是 Hausdorff 的，当且仅当它是 P'-无挠的。
\end{enumerate}

\begin{enumerate}
\def\labelenumi{\arabic{enumi}.}
\setcounter{enumi}{17}
\tightlist
\item
  群 G 的上中心列 \((Z_{i}G)\) 递归定义如下：
\end{enumerate}

\begin{enumerate}
\def\labelenumi{(\roman{enumi})}
\item
  \(Z_{i}G = \{e\}\) 当 \(i \leqslant 0\) 时；
\item
  \(Z_{i}G/Z_{i-1}G\) 是 \(G/Z_{i-1}G\) 的中心。
\end{enumerate}

于是 \(\{e\} = Z_0G \subset Z_1G \subset \cdots\)，且 \(Z_1G\) 是 G 的中心；\(Z_iG\) 是 G 的特征子群。

\begin{enumerate}
\def\labelenumi{(\alph{enumi})}
\item
  证明 \(\mathbf{G}\) 是类 \(\leqslant c\) 的幂零群，当且仅当 \(\mathbf{G} = \mathbf{Z}_c\mathbf{G}\)。
\item
  证明 \((\mathbf{C}^n\mathbf{G},\bar{\mathbf{Z}}_m\mathbf{G})\subset \mathbf{Z}_{m - n}\mathbf{G}.\)
\item
  假设 G 是幂零且 P-无挠的（P 是素数集合，参见~习题 14）。证明 \(Z_{1}G\) 是 P-饱和的。（只需对 \(Z_{1}G\) 验证；若 n 是 P-整数，且 \(g \in G\) 满足 \(g^{n} \in Z_{1}G\)，则 \(xg^{n}x^{-1} = g^{n}\) 对所有 \(x \in G\) 成立，由习题 14 (c) 得 \(xgx^{-1} = g\)，从而 g 属于 \(Z_{1}G\)。）
\end{enumerate}

\exercisesfor{5}

在下列习题中，假设 §5 的假设和记号成立。F 表示自由群 \(\mathrm{F}(\mathbf{X})\)，g 表示从 F 到 Magnus 群 \(\Gamma(\mathbf{X})\) 的唯一同态，满足 \(g(x)=1+x\) 对所有 \(x\in X\) 成立（参见~定理 1）。

\begin{enumerate}
\def\labelenumi{\arabic{enumi}.}
\tightlist
\item
  给代数 \(\mathbf{K}[\mathbf{F}]\)（\(\mathbf{F}\) 的代数）赋予由 \((\mathbf{I}^n)\)（增广理想 \(\mathbf{I}\) 的各次幂）组成的滤过（\(\S 4\)，习题 6）。
\end{enumerate}

\begin{enumerate}
\def\labelenumi{(\alph{enumi})}
\item
  令 \(\tilde{g}\colon \mathrm{K}[\mathrm{F}] \to \tilde{\mathrm{A}}(\mathrm{X})\) 为延拓 \(g\colon \mathrm{F} \to \tilde{\mathrm{A}}(\mathrm{X})^*\) 的唯一代数同态。证明 \(\tilde{g}\) 将 \(\mathrm{I}^{\mathfrak{n}}\) 映到理想 \(\tilde{\mathrm{A}}_{\mathfrak{n}}(\mathrm{X})\) 中（参见~第 1 号），并且在取商后定义同构 \(\tilde{g}_{n}\)，其从 \(\mathrm{K}[\mathrm{F}]/\mathrm{I}^{\mathfrak{n}}\) 到 \(\tilde{\mathrm{A}}(\mathrm{X})/\tilde{\mathrm{A}}_{\mathfrak{n}}(\mathrm{X})\)。（利用 \(\tilde{g}_{n}\) 的逆同态，借助从 \(\mathrm{A}(\mathrm{X})\) 到 \(\mathrm{K}[\mathrm{F}]\) 的同态，该同态将 \(x\) 映为 \(x - 1\)，对所有 \(x \in \mathrm{X}\) 成立。）推出 \(\tilde{\mathrm{A}}(\mathrm{X})\) 同构于 \(\mathrm{K}[\mathrm{F}]\) 关于由 \((\mathrm{I}^{\mathfrak{n}})\) 定义的拓扑的豪斯多夫完备化。
\item
  假设 \(\mathbf{K} = \mathbf{Z}\)。证明滤过 \((\mathrm{I}^{\mathrm{n}})\) 是分离的（使用命题 3 和 §4 的习题 6 (e)），并且由它定义的 \(\mathbf{F}\) 的滤过（§4 的习题 6 (a)）与滤过 \((\mathbf{C}^{\mathrm{n}}\mathbf{F})\) 重合。推出 \(\tilde{g}\colon \mathbf{Z}[\mathrm{F}] \to \hat{\mathrm{A}}_{\mathbf{Z}}(\mathrm{X})\) 是单射。
\end{enumerate}

\begin{enumerate}
\def\labelenumi{\arabic{enumi}.}
\setcounter{enumi}{1}
\tightlist
\item
  设 G 是一个群，K{[}G{]} 是它在 K 上的代数，I 是其增广理想（§4，习题 6）。令 \(\varepsilon\) 表示从 K{[}G{]} 到 K 上的典范同态；于是 \(\operatorname{Ker}(\varepsilon) = I\)。
\end{enumerate}

\begin{enumerate}
\def\labelenumi{(\alph{enumi})}
\tightlist
\item
  设 M 是左 K{[}G{]}-模，Z(G, M) 是从 G 到 M 的交叉同态群（代数，第一章，§6，习题 7）。若 \(f \in \operatorname{Hom}_{\mathbb{K}[\mathbb{G}]}(\mathrm{I}, \mathrm{M})\)，令 \(f_1\) 为映射 \(g \mapsto f(g - 1)\)（从 \(\mathbf{G}\) 到 \(\mathbf{M}\)）。证明 \(f_1\) 是交叉同态（利用恒等式
\end{enumerate}

\[
g g ^ {\prime} - 1 = g (g ^ {\prime} - 1) + g - 1
\]

证明 \(f_{1}(gg') = g \cdot f_{1}(g') + f_{1}(g)\)，并证明 \(f \mapsto f_{1}\) 是从 \(\operatorname{Hom}_{\mathbb{K}[G]}(\mathbf{I}, \mathbf{M})\) 到 \(\mathbf{Z}(\mathbf{G}, \mathbf{M})\) 的同构。

\begin{enumerate}
\def\labelenumi{(\alph{enumi})}
\setcounter{enumi}{1}
\item
  取 G 为自由群 \(\mathrm{F} = \mathrm{F}(\mathrm{X})\)。证明对每个映射 \(\theta: \mathrm{X} \to \mathrm{M}\)，存在唯一的元素 \(f_{\theta}\)，属于 \(\mathrm{Z}(\mathrm{M}, \mathrm{F})\) 且延拓 \(\theta\)。（利用将 \(\mathrm{Z}(\mathrm{F}, \mathrm{M})\) 解释为 F 与 M 的半直积，参见~代数第一章上述位置。）利用 (a) 推出，族 \((x - 1)_{x \in \mathbb{X}}\) 是 I 作为左 K{[}F{]}-模的一组基。
\item
  由 (b)，每个元素 \(u \in \mathbf{K}[\mathbf{F}]\) 都能唯一写成如下形式：
\end{enumerate}

\[
u = \varepsilon (u) + \sum_ {x \in \mathbf {X}} \mathrm{D} _ {x} (u) (x - 1),
\]

其中 \(\mathrm{D}_{x}(u)\in\mathrm{K}[\mathrm{F}]\)，且除有限多个外，\(\mathrm{D}_{x}(u)\) 均为零。映射 \(u\mapsto\mathrm{D}_{x}(u)\) 称为关于 x 的偏导数。它由下列性质刻画：

\begin{enumerate}
\def\labelenumi{(\roman{enumi})}
\item
  \(D_{x}\) 是从 K{[}F{]} 到 K{[}F{]} 的 K-线性映射。
\item
  \(\mathrm{D}_{x}(uv)=u.\mathrm{D}_{x}(v)+\varepsilon(v).\mathrm{D}_{x}(u)\)，其中 u、v 属于 K{[}F{]}。
\item
  \(\mathrm{D}_{x}(x)=1\)，且 \(\mathrm{D}_{x}(y)=0\) 当 \(y\in X-\{x\}\) 时。
\end{enumerate}

若 \(n\geqslant 1\)，则

\[
\begin{array}{c} \mathrm{D} _ {x} (x ^ {n}) = 1 + x + \dots + x ^ {n - 1} \\ \mathrm{D} _ {x} (x ^ {- n}) = - x ^ {- n} - x ^ {- n + 1} - \dots - x ^ {- 1}. \end{array}
\]

\begin{enumerate}
\def\labelenumi{(\alph{enumi})}
\setcounter{enumi}{3}
\tightlist
\item
  令 \(\varepsilon_{\mathrm{A}}\) 为从 A(X) 到 K 的典范同态。证明 A(X) 的每个元素 \(u\) 都能唯一写成如下形式：
\end{enumerate}

\[
u = \varepsilon_ {\mathrm{A}} (u) + \sum_ {x \in \mathbf {X}} d _ {x} (u). x,
\]

其中 \(d_{x}(u) \in \mathrm{A}(\mathrm{X})\)，且除有限多个外，\(d_{x}(u)\) 均为零。证明 \(d_{x}: u \to d_{x}(u)\) 是 \(\mathrm{A}(\mathrm{X})\) 的自同态，可以通过连续性延拓到 \(\hat{\mathrm{A}}(\mathrm{X})\)，并具有类似于上面 (i)、(ii)、(iii) 的性质。证明 \(d_{x} \circ \tilde{g} = \tilde{g} \circ D_{x}\)，其中 \(\tilde{g}\) 是延拓 g 的从 K{[}F{]} 到 \(\hat{\mathrm{A}}(\mathrm{X})\) 的同态（参见~习题 1）。†

¶ 3. 沿用习题 2 的记号。

\begin{enumerate}
\def\labelenumi{(\alph{enumi})}
\tightlist
\item
  设 \(J\) 是 \(K[F]\) 的右理想，且包含于 \(I\)，并设 \(u\) 是 \(I\) 的一个元素。证明
\end{enumerate}

\[
u \in J. I \Leftrightarrow D _ {x} (u) \in J \quad \text { for   all } x \in X.
\]

特别地，I 的元素属于 \(\mathbf{I}^n\)（\(n \geqslant 1\)），当且仅当它的所有偏导数都属于 \(\mathbf{I}^{n-1}\)。


\begin{enumerate}
\def\labelenumi{(\alph{enumi})}
\setcounter{enumi}{1}
\tightlist
\item
  令 R 是 F 的正规子群，令 G = F/R，令 J 是典范同态 \(\gamma: K[F] \to K[G]\) 的核。证明 J 作为左理想（相应地，右理想）由元素 r - 1（其中 \(r \in R\)）生成。证明下列序列正合：
\end{enumerate}

\[
0 \longrightarrow \mathrm{J} / (\mathrm{J}. \mathrm{I}) \stackrel {\alpha} {\longrightarrow} \mathrm{I} / (\mathrm{J}. \mathrm{I}) \stackrel {\beta} {\longrightarrow} \mathrm{K} [ \mathrm{G} ] \stackrel {\delta} {\longrightarrow} \mathrm{K} \longrightarrow 0,\tag{*}
\]

其中 \(\alpha\) 由 J 包含于 I 的包含映射取商得到，\(\beta\) 由 \(\gamma\) 在 I 上的限制取商得到。

\begin{enumerate}
\def\labelenumi{(\alph{enumi})}
\setcounter{enumi}{2}
\item
  对 \(u \in \mathbf{I}\)、\(x \in \mathbf{X}\)，令 \(\overline{\mathrm{D}}_x(u)\) 表示 \(\mathrm{D}_x(u)\) 在 \(\mathrm{K}[\mathrm{G}]\) 中的像（由 \(\gamma\) 给出）。利用 (a) 证明，族 \((\overline{\mathrm{D}}_x)_{x \in \mathbf{X}}\) 在取商后定义了从 \(\mathrm{I} / (\mathrm{J}. \mathrm{I})\) 到 \(\mathrm{K}[\mathrm{G}]^{(\mathbb{X})}\) 的同构。
\item
  对 \(r \in \mathbb{R}\)，令 \(\theta(r)\) 表示 \(r - 1\) 在 \(J/(J.I)\) 中的像。证明
\end{enumerate}

\[
\theta (r r ^ {\prime}) = \theta (r) + \theta (r ^ {\prime}) \quad \text { and } \quad \theta (y r y ^ {- 1}) = y. \theta (r) \quad \text { for   } r, r ^ {\prime} \text {   in   } R, y \in F.
\]

（使用恒等式

\[
\begin{array}{c} {r r ^ {\prime} - 1 = (r - 1) (r ^ {\prime} - 1) + (r - 1) + (r ^ {\prime} - 1)} \\ {y r y ^ {- 1} - 1 = y (r - 1) (y ^ {- 1} - 1) + y (r - 1).)} \end{array}
\]

推出 \(\theta\) 定义了一个从 \(R/(R, R)\) 到 \(J/(J, I)\) 的同态，并且与 \(G\) 的作用相容；证明该同态的像生成 K-模 \(J/(J, I)\)。当 \(K = Z\) 时，证明由此得到从 \(R/(R, R)\) 到 \(J/(J, I)\) 的同构（直接定义逆同态）。（e）令 \((r_{\alpha})_{\alpha \in A}\) 是 \(R\) 中的一个族，且它生成 \(R\) 作为 \(F\) 的正规子群。证明 \(\theta(r_{\alpha})\) 生成 K{[}G{]}-模 \(J/(J, I)\)。

\begin{enumerate}
\def\labelenumi{(\alph{enumi})}
\setcounter{enumi}{5}
\tightlist
\item
  矩阵 \((\overline{\mathrm{D}}_x(r_\alpha))_{x\in \mathbf{X},\alpha \in \mathbf{A}}\) 定义了一个同态
\end{enumerate}

\[
\rho \colon \mathrm{K} [ \mathrm{G} ] ^ {(\mathrm{A})} \rightarrow \mathrm{K} [ \mathrm{G} ] ^ {(\mathrm{X})}
\]

它是左 K{[}G{]}-模的同态。证明下列序列正合：

\[
\mathrm{K} [ \mathrm{G} ] ^ {(\mathrm{A})} \xrightarrow {\rho} \mathrm{K} [ \mathrm{G} ] ^ {(\mathrm{X})} \xrightarrow {\delta} \mathrm{K} [ \mathrm{G} ] \xrightarrow {\varepsilon} \mathrm{K} \longrightarrow 0\tag{**}
\]

其中 \(\delta\) 是同态 \((u_x) \mapsto \sum_{x \in X} u_x(\gamma(x) - 1)\)。

（利用上面的 (c)、(d)、(e) 变换正合序列 (*)。）†

\begin{enumerate}
\def\labelenumi{\arabic{enumi}.}
\setcounter{enumi}{3}
\tightlist
\item
  对所有 \(n \in \mathbf{N}\)，令 \(\binom{\mathrm{T}}{n}\) 表示多项式
\end{enumerate}

\[
\mathrm{T} (\mathrm{T} - 1) \dots (\mathrm{T} - n + 1) / n!.
\]

若 \(f(\mathrm{T})\) 是多项式，令 \(\Delta f\) 表示多项式 \(f(\mathrm{T} + 1) - f(\mathrm{T})\)。

于是 \(\Delta\binom{\mathrm{T}}{0} = \Delta 1 = 0\)，并且 \(\Delta\binom{\mathrm{T}}{n} = \binom{\mathrm{T}}{n - 1}\) 在 \(n \geqslant 1\) 时成立。

\begin{enumerate}
\def\labelenumi{(\alph{enumi})}
\item
  令 \(f \in \mathbf{Q}[\mathrm{T}]\)。证明以下性质等价：
\item
  \(f\) 将 \(\mathbf{Z}\) 映入自身。
\end{enumerate}

\begin{enumerate}
\def\labelenumi{(\roman{enumi})}
\setcounter{enumi}{1}
\item
  \(f\) 是系数属于 \(\mathbf{Z}\) 的各个 \(\binom{\mathrm{T}}{n}\) 的线性组合。
\item
  \(f(0)\in \mathbf{Z}\)，且 \(\Delta f\) 将 \(\mathbf{Z}\) 映入自身。
\end{enumerate}

（对 \(\deg(f)\) 归纳，并注意 \(\deg(\Delta f) = \deg(f) - 1\) 在 \(\deg(f) \neq 0\) 时成立。）

满足上述性质的多项式 \(f\) 称为二项式多项式。两个二项式多项式的和、积和复合仍是二项式多项式。

\begin{enumerate}
\def\labelenumi{(\alph{enumi})}
\setcounter{enumi}{1}
\item
  令 \(p\) 是素数，并令 \(f\) 是二项式多项式。证明 \(f\) 将环 \(\mathbf{Z}_p\)（\(p\)-进整数环）映入自身（利用 \(f\) 的连续性以及 \(\mathbf{Z}\) 在 \(\mathbf{Z}_p\) 中稠密这一事实）。
\item
  令 \(\mathbf{P}\) 是素数集合，\(\mathrm{S}_P\) 是 P-整数的集合（§4，习题 14），令 \(\mathbf{Z}_P = \mathrm{S}_{\mathrm{P}}^{-1}\mathbf{Z}\)。证明若 \(f\) 是二项式多项式，则 \(f\) 将 \(\mathbf{Z}_{\mathrm{P}}\) 映入自身（将 (b) 应用于 \(p\) 不属于 \(\mathbf{P}\) 的素数）。
\end{enumerate}

\begin{enumerate}
\def\labelenumi{\arabic{enumi}.}
\setcounter{enumi}{4}
\tightlist
\item
  令 \(\Gamma = \Gamma(X)\) 是 \(X\) 上的 Magnus 群，令 \((\Gamma_n)\) 是其自然滤过（第 2 号）。证明与此滤过相应的分次李代数 \(\mathrm{gr}(\Gamma)\) 同构于李子代数 \(\bigoplus_{n\geqslant 1}A_n(X)\)，它是 \(A(X)\) 的子代数。若 \(X\neq \emptyset\)，此李代数不由其次数为 1 的元素生成；推出 \((\Gamma_n)\) 不是 \(\Gamma\) 的下中心列。
\end{enumerate}

¶ 6. 令 P 是素数集合，令 \(S_P\) 为 P-整数的集合（§4，习题 14），令 \(Z_P = S_P^{-1}Z\)。

\begin{enumerate}
\def\labelenumi{(\alph{enumi})}
\tightlist
\item
  假设对所有 \(s \in S_{P}\)，映射 \(k \mapsto sk\) 都是 K 到自身的双射；这等价于说 K 可以赋予 \(Z_{P}\)-代数结构。
\end{enumerate}

沿用习题 5 的记号，证明对所有 \(s \in S_{\mathbb{P}}\)，映射 \(a \mapsto a^s\) 将 \(\Gamma\) 双射到自身，并且商群 \(\Gamma / \Gamma_n\) 对每个 \(n \geqslant 1\) 也具有同样性质。若 \(t \in \mathbf{Z}_{\mathbb{P}}\) 且 \(a \in \Gamma\)（相应地，\(a \in \Gamma / \Gamma_n\)），则按 §4 习题 16 定义 \(a^s\)；将 \(a\) 写成 \(1 + \alpha\) 的形式，其中 \(\alpha \in \hat{\mathrm{A}}_1(\mathrm{X})\)，证明

\[
a ^ {t} = (1 + \alpha) ^ {t} = \sum_ {n = 0} ^ {\infty} \binom {t} {n} \alpha^ {n}.
\]

（注意系数 \(\binom{t}{n}\) 属于 \(\mathbf{Z}_{\mathrm{P}}\)，参见~习题 4。）

\begin{enumerate}
\def\labelenumi{(\alph{enumi})}
\setcounter{enumi}{1}
\item
  现在假设 \(\mathbf{K} = \mathbf{Z}_{\mathrm{P}}\)。令 \(c\) 是满足 \(\geqslant 1\) 的整数。将群 \(\mathbf{F}^c = \mathbf{F} / \mathbf{C}^c\mathbf{F}\) 与 \(\Gamma /\Gamma_c\) 的一个子群相识别，所用的是由 \(g\) 取商得到的同态（定理 2 的记号）。群 \(\Gamma / \Gamma_c\) 是类 \(< c\) 的 P-无挠、P-可除群。令 \(F_P^c\) 为 \(F^c\) 在 \(\Gamma / \Gamma_c\) 中的 P-饱和（§4，习题 14），并令 \(i\) 是从 \(F^c\) 到 \(F_P^c\) 的嵌入。序对 \((i, F_P^c)\) 是 \(F^c\) 的 P-包络（§4，习题 15）。推出每个幂零群都有 P-包络（注意每个类 \(< c\) 的幂零群都是某个适当 X 的群 \(F^c\) 的商，并使用§4 习题 15 的 (b) 部分）。
\item
  若 \(n \leqslant c\)，令 \(\mathbf{F}_{\mathbf{P},n}^{\mathrm{c}}\) 为 \(\mathbf{F}_{\mathbf{P}}^{\mathrm{c}}\) 与 \(\Gamma_{n}/\Gamma_{c}\) 的交；若 \(n \geqslant c\)，记 \(\mathbf{F}_{\mathbf{P},n}^{\mathrm{c}} = \{e\}\)。证明 \(\mathbf{F}_{\mathbf{P},n}^{\mathrm{c}}\) 是 \(\mathbf{C}^{\mathrm{c}}\mathbf{F}^{\mathrm{c}}\) 在 \(\Gamma/\Gamma_{c}\) 中的 P-饱和。滤过 \((\mathbf{F}_{\mathbf{P},n}^{\mathrm{c}})\) 是 \(\mathbf{F}_{\mathbf{P},n}^{\mathrm{c}}\) 的整中心滤过；令 \(\text{gr}(\mathbf{F}_{\mathbf{P}}^{\mathrm{c}})\) 为与此滤过相伴的分次模。证明若 \(n < c\)，则 \(\text{gr}_{\mathbf{n}}(\mathbf{F}_{\mathbf{P}}^{\mathrm{c}})\) 在 \(\text{gr}_{\mathbf{n}}(\Gamma) = \mathbf{A}_{\mathbf{Z}_{\mathbf{P}}}^{n}(\mathbf{X})\) 中的像是 \(\mathbf{S}_{\mathbf{P}}^{-1}. \mathbf{L}_{\mathbf{Z}}^{n}(\mathbf{X}) = \mathbf{L}_{\mathbf{Z}_{\mathbf{P}}}^{n}(\mathbf{X})\)。推出李代数 \(\text{gr}(\mathbf{F}_{\mathbf{P},n}^{\mathrm{c}})\) 由其次数为 1 的元素生成，从而 \(\mathbf{F}_{\mathbf{P},n} = \mathbf{C}^{n}(\mathbf{F}_{\mathbf{P}}^{\mathrm{c}})\) 对所有 \(n\) 都成立（\(\S 4\)，习题 1）。证明群 \(\mathbf{F}_{\mathbf{P}}^{\mathrm{c}}\) 由 \(x^{1/s}\)（其中 \(x \in \mathbf{X}\)、\(s \in \mathbb{S}_{\mathbf{P}}\)）生成（注意这些元素在 \(\text{gr}_1(\mathbf{F}_{\mathbf{P}}^{\mathrm{c}}) \simeq \mathbf{Z}_{\mathbf{P}}^{(X)}\) 中的像生成群 \(\text{gr}_1(\mathbf{F}_{\mathbf{P}}^{\mathrm{c}})\)，并应用 \(Algebra\) 第一章 §6 第 3 号命题 8 的推论 3）。
\item
  令 \(\mathrm{H}\) 是相对于 \(\mathbf{X}\) 的 Hall 集（§2，第 10 号），令 \(\mathrm{H}(c)\) 为由 \(\mathrm{H}\) 中长度 \(< c\) 的元素组成的子集。对 \(m \in \mathrm{H}(c)\)，令 \(\phi_c(m)\) 表示基本换位子在 \(\mathbf{F}^c\) 中的像；这里的基本换位子是 \(\phi(m)\)，由 \(m\) 定义（参见~第 4 号的注）。证明对所有 \(w \in \mathbf{F}_P^c\)，存在唯一的 \(\alpha \in \mathbf{Z}_{\mathbf{P}}^{\mathbf{H}(c)}\)，使得
\end{enumerate}

\[
w = \prod_ {m \in H (c)} \phi_ {c} (m) ^ {\alpha (m)}.
\]

（使用上面得到的 gr(F \(_{P}^{c}\) ) 的刻画。）

\begin{enumerate}
\def\labelenumi{\arabic{enumi}.}
\setcounter{enumi}{6}
\tightlist
\item
  沿用习题 6 的记号。
\end{enumerate}

\begin{enumerate}
\def\labelenumi{(\alph{enumi})}
\item
  令 G 是类 \(<c\) 的幂零群，令 \((i, \overline{G})\) 是 G 的一个 P-包络。令 \(\rho: \mathrm{F}^{c} \to \mathrm{G}\) 是满同态（适当选取 X 时这样的同态存在），令 \(\bar{\rho}\) 是相应的 \(\mathrm{F}_{\mathbf{P}}^{c} \in \overline{G}\) 的同态（§4，习题 15）。同态 \(\bar{\rho}\) 是满射，并将 \(\mathrm{C}^{n}\mathrm{F}_{\mathbf{P}}^{c} = \mathrm{F}_{\mathbf{P}, n}^{c}\) 映到 \(\mathrm{C}^{n}\overline{G}\) 上。证明 \(\mathrm{C}^{n}\overline{G}\) 是 \(i(\mathrm{C}^{n}\mathrm{G})\) 在 \(\overline{G}\) 中的 P-包络，并且李代数 \(\mathrm{gr}(\overline{G})\) 可与 \(\mathbf{Z}_{\mathbf{P}} \otimes \mathrm{gr}(\mathrm{G})\) 相识别。
\item
  推出 G 是 P-可除且 P-无挠的，当且仅当各个 \(C^{n}G/C^{n+1}G\) 具有这些性质。
\end{enumerate}

¶ 8. 设 K = Z 且 X 有限。对每个整数 \(k \geqslant 0\)，令 \((e_{k,\alpha})\) 表示 \(A_k(X)\) 的一组基。

\begin{enumerate}
\def\labelenumi{(\alph{enumi})}
\tightlist
\item
  令 \((w_{n})_{n\in \mathbf{N}}\) 是 F 中元素的一个序列。令 \(w_{k,\alpha}(n)\) 为 \(e_{k,\alpha}\) 的系数；该系数取自 \(g(w_{n})\in \hat{\mathrm{A}} (\mathrm{X})\) 中次数为 \(k\) 的项。则：
\end{enumerate}

\[
g (w _ {n}) = \sum_ {k, \alpha} w _ {k, \alpha} (n) e _ {k, \alpha} \quad \text { for   all } n \in \mathbf {Z}.
\]

序列 \((w_{n})\) 称为典型的，如果对每个有序对 \((k,\alpha)\)，函数 \(w_{k,\alpha}\colon n\mapsto w_{k,\alpha}(n)\) 是次数 \(\leqslant k\) 的二项式多项式，参见~习题 4。此条件与基 \((e_{k,\alpha})\) 的选择无关。证明它等价于存在一个序列 \((a_{k})_{k\in \mathbf{N}}\)，其元素属于 \(\hat{\mathbf{A}}_{\mathbf{q}}(\mathbf{X})\)，使得 \(\omega (a_k)\geqslant k\) 对所有 \(k\) 都成立，并且

\[
g (w _ {n}) = \sum_ {k = 0} ^ {\infty} n ^ {k} a _ {k} \quad \text { for   all } n \in \mathbf {Z}.
\]

\begin{enumerate}
\def\labelenumi{(\alph{enumi})}
\setcounter{enumi}{1}
\item
  证明若 \((w_{n})\) 和 \((w_{n}^{\prime})\) 是两个典型序列，则 \((w_{n}^{-1})\) 和 \((w_{n}w_{n}^{\prime})\) 也都是典型序列。
\item
  令 \(w\) 是 \(\mathbf{C}^k\mathbf{F}\) 中的元素，其中 \(k \geqslant 1\)，令 \(f\) 是次数 \(\leqslant k\) 的二项式多项式。证明 \((w^{f(n)})\) 是典型序列。特别地，若 \(z \in \mathbf{F}\)，则幂序列 \((z^n)\) 是 \(z\) 的典型幂序列。
\item
  令 \((w_{n})\) 是 F 中元素的一个序列。证明 \((w_{n})\) 是典型的，当且仅当存在 F 中的 \(y_{0}, y_{1}, \ldots, y_{n}, \ldots\)，满足 \(y_{i} \in \mathbf{C}^{i}\mathbf{F}\) 对所有 \(i \geqslant 1\) 都成立，并且
\end{enumerate}

\[
w _ {n} \equiv y _ {0} y _ {1} ^ {n} \dots y _ {k} ^ {\binom {n} {k}} \mod \mathrm{C} ^ {k + 1} \mathrm{F}
\]

对所有 \(n \in \mathbf{Z}\) 和所有 \(k \geqslant 1\) 均成立。（确定 \(y_{i}\)：由 \((w_{n})\) 依次令 \(n = 0, 1, \ldots\)。例如 \(w_{0} = y_{0}, w_{1} = y_{0}y_{1}, \ldots\)。）

\begin{enumerate}
\def\labelenumi{(\alph{enumi})}
\setcounter{enumi}{4}
\tightlist
\item
  令 H 是相对于 X 的 Hall 集，\((w_{n})\) 是 F 中元素的一个序列。令 \((\alpha(m,n))_{m\in\mathbb{H},n\in\mathbf{Z}}\) 为满足下式的整数族：
\end{enumerate}

\[
w _ {n} \equiv \prod_ {m \in \mathrm{H}} \phi (m) ^ {\alpha (m, n)} \mod \mathrm{C} ^ {c} \mathrm{F}
\]

对所有 \(n \in \mathbf{Z}\) 和所有 \(c \geqslant 1\) 均成立（参见~第 4 号的注）。证明 \((w_n)\) 是典型的，当且仅当对所有 \(m \in \mathbf{H}\)，函数 \(n \mapsto \alpha(m, n)\) 是次数 \(\leqslant l(m)\) 的二项式多项式，其中 \(l(m)\) 是 \(m\) 的长度。†

\begin{enumerate}
\def\labelenumi{(\alph{enumi})}
\setcounter{enumi}{5}
\tightlist
\item
  序列 \((w_{n})\) 称为 1-典型，如果相应的函数 \(w_{k,\alpha}\) 是次数 \(\leqslant k - 1\) 的二项式多项式。证明若 \((w_{n})\) 是 1-典型的，则存在 \(y_{i} \in \mathbf{C}^{i + 1}\mathbf{F}\)（\(i = 0, 1, \ldots\)），使得
\end{enumerate}

\[
w _ {n} \equiv y _ {0} y _ {1} ^ {n} \dots y _ {k} ^ {\binom {n} {k}} \mod \mathrm{C} ^ {k + 1} \mathrm{F}
\]

对所有 \(n \in Z\) 和所有 \(k \geqslant 1\) 均成立。

¶ 9. 令 X 是含两个元素 \{x, y\} 的集合。

\begin{enumerate}
\def\labelenumi{(\alph{enumi})}
\tightlist
\item
  证明存在一个序列 \(w_{2}, w_{3}, \ldots\)，其元素属于 \(\mathbf{F}\)，且 \(w_{i} \in \mathbb{C}^{i}\mathbb{F}\) 对所有 \(i\) 都成立，并且
\end{enumerate}

\[
(x y) ^ {n} \equiv x ^ {n} y ^ {n} w _ {2} ^ {\binom {n} {2}} \dots w _ {i} ^ {\binom {n} {i}} \mod \mathrm{C} ^ {i + 1} \mathrm{F}
\]

对所有 \(n \in \mathbf{Z}\) 和所有 \(i \geqslant 1\) 均成立。（将习题 8（\(d\)）应用于序列 \(x^{-n}(xy)^n\)。）证明 \(w_2 = y^{-1}(y^{-1}, x^{-1})y \equiv (x, y)^{-1} \mod \mathrm{C}^3\mathrm{F}\)。

\begin{enumerate}
\def\labelenumi{(\alph{enumi})}
\setcounter{enumi}{1}
\tightlist
\item
  设 G 是类 \(\leqslant c\) 的幂零群，且 x、y 属于 G。由上文推出下列公式（“Hall 公式”）：
\end{enumerate}

\[
(x y) ^ {n} = x ^ {n} y ^ {n} \prod_ {i = 2} ^ {c} w _ {i} (x, y) ^ {\binom {n} {i}}.
\]

\begin{enumerate}
\def\labelenumi{(\alph{enumi})}
\setcounter{enumi}{2}
\tightlist
\item
  令 \(p\) 为素数。证明存在 \(u_{i} \in \mathbf{C}^{i}\mathbf{F}\)（ \(2 \leqslant i \leqslant p - 1\) ）以及 \(w \in \mathbf{C}^{p}\mathbf{F}\)，使得
\end{enumerate}

\[
(x y) ^ {p} = x ^ {p} y ^ {p} u _ {2} ^ {p} \dots u _ {p - 1} ^ {p} w.
\]

（使用 (a)，注意 \(\binom{p}{i}\) 被 \(p\) 整除，只要 \(1 < i < p\)。）令 \(\bar{w}\) 为 \(w\) 在 \(\mathrm{gr}_p(\mathsf{F}) = \mathrm{L}_{\mathbf{Z}}^p (\mathbf{X})\) 中的像，令 \(\bar{w}_p\) 为 \(\bar{w}\) 在 \(\mathrm{L}_{\mathbf{F}_p}^p (\mathbf{X})\) 中的像。证明 \(\bar{w}_p = \Lambda_p(x,y)\)，参见~第一章，§ 1，习题 19。（将 F 通过同态 \(g\) 延拓到 \(\tilde{\mathbf{A}}(\mathbf{X})\)，并比较次数为 \(p\) 的 \(g((xy)^p)\) 与 \(g(x^p y^p u_2^p \ldots u_{p-1}^p w)\) 中的项。）前者等于 \((x + y)^p\)，后者模 \(p\) 同余于 \(x^p + y^p + \bar{w}_p\)。故结论成立。）

证明存在 \(a \in \mathbf{C}^2\mathbf{F}\) 以及 \(w' \in \mathbf{C}^p\mathbf{F}\)，使得

\[
(x y a) ^ {p} = x ^ {p} y ^ {p} w ^ {\prime}.
\]

（使用上式计算 \((xy)^p\)。）由此推出，类为 \(p\) 的幂零群中，\(p\) 次幂构成一个子群。

\begin{enumerate}
\def\labelenumi{(\alph{enumi})}
\setcounter{enumi}{3}
\tightlist
\item
  证明存在一个序列 \(v_{3}, v_{4}, \ldots\)，其元素属于 \(\mathbf{F}\)，且 \(v_{i} \in \mathbf{C}^{i}\mathbf{F}\) 对所有 \(i\) 都成立，并且
\end{enumerate}

\[
(x ^ {n}, y) \equiv (x, y) ^ {n} v _ {3} ^ {\binom {n} {2}} \dots v _ {i + 1} ^ {\binom {n} {i}} \mod \mathrm{C} ^ {i + 2} \mathrm{F}
\]

对所有 \(n \in \mathbf{N}\) 和所有 \(i \geqslant 1\) 均成立。（将习题 8(f) 应用于 \((x^n, y)\) 的序列。）

推出，若 \(p\) 为素数，则存在 \(t_i \in \mathrm{C}^i\mathrm{F}\)（ \(3 \leqslant i \leqslant p\)）以及 \(z \in \mathrm{C}^{p+1}\mathrm{F}\)，使得

\[
(x ^ {p}, y) = (x, y) ^ {p} t _ {3} ^ {p} \dots t _ {p} ^ {p} z.
\]

证明 \(\bar{z}_p\) 是 \(z\) 在 \(\mathrm{L}_{\mathbf{F}_p}^{p + 1}(\mathbf{X})\) 中的像，且等于 (ad \(x)^p (y)\)。（方法同 \((c).)\)

¶ 10. 令 p 为素数，令 G 是带有实中心滤过 (G \(_{\alpha}\) ) 的群。若关系 x ∈ G \(_{\alpha}\) 蕴含 x \(^{p}\) ∈ G \(_{p\alpha}\)，则称滤过 (G \(_{\alpha}\) ) 为受限滤过。

\begin{enumerate}
\def\labelenumi{(\alph{enumi})}
\item
  证明李代数 \(\operatorname{gr}(\mathbf{G})\)（与 \((\mathbf{G}_{\alpha})\) 相关）满足 \(p.\operatorname{gr}(\mathbf{G}) = 0\)，因而可以赋予 \(\mathbf{F}_p\) 上的代数结构。
\item
  令 \(\xi \in \mathrm{gr}_{\alpha}(\mathbf{G})\)，令 \(x\) 是 \(\xi\) 在 \(\mathbf{G}_{\alpha}\) 中的一个代表元。证明 \(x^{p}\) 在 \(\mathrm{gr}_{p\alpha}(\mathbf{G})\) 中的像与 \(x\) 的选取无关（使用习题 9(c)）。若将此像记为 \(\xi^{[p]}\)，证明 \(\xi \mapsto \xi^{[p]}\) 是从 \(\mathrm{gr}_{\alpha}(\mathbf{G})\) 到 \(\mathrm{gr}_{p\alpha}(\mathbf{G}')\) 的线性映射，并且 \((\xi + \xi')^{[p]} = \xi^{[p]} + \xi'^{[p]} + \Delta_p(\xi, \xi')\)。（方法相同。）
\end{enumerate}

证明，若 \(\xi \in \mathrm{gr}_{\alpha}(\mathbf{G})\) 且 \(\eta \in \mathrm{gr}_{\beta}(\mathbf{G})\)，则

\[
[ \xi^ {[ p ]}, \eta ] = (\operatorname{ad} \xi) ^ {p} (\eta).
\]

（使用习题 9(d)。）

\begin{enumerate}
\def\labelenumi{(\alph{enumi})}
\setcounter{enumi}{2}
\tightlist
\item
  证明存在唯一的 \(p\)-映射，它是从 \(\operatorname{gr}(\mathbf{G})\) 到自身的映射（第一章，§ 1，习题 20），且延拓了映射 \(\xi \mapsto \xi^{[p]}\)，该映射定义在各个 \(\operatorname{gr}_{\alpha}(\mathbf{G})\) 上。赋予这一结构后，\(\operatorname{gr}(\mathbf{G})\) 是一个李 \(p\)-代数，称为与受限滤过 \((\mathbf{G}_{\alpha})\) 相关的李代数。
\end{enumerate}

¶ 11. (a) 令 A 是满足 §4 第 5 号条件的滤过代数，令 \(\Gamma = A^{*} \cap (1 + A_{0}^{+})\)。在 A 上的滤过诱导出 \(\Gamma\) 上的滤过（同上，命题 2）。证明若 K 是特征为 \(p > 0\) 的域，则 \((\Gamma_{\alpha})\) 是受限滤过（习题 10），并且同上命题 3 所定义的 \(\mathrm{gr}(\Gamma)\) 到 \(\mathrm{gr(A)}\) 的嵌入与李 \(p\)-代数结构相容；这些结构分别定义在 \(\mathrm{gr}(\Gamma)\) 和 \(\mathrm{gr(A)}\) 上。

\begin{enumerate}
\def\labelenumi{(\alph{enumi})}
\setcounter{enumi}{1}
\item
  假设 \(\mathbf{K} = \mathbf{F}_p\)。将上述结果应用于滤过代数 \(\hat{\mathbf{A}}(\mathbf{K})\)，其中 \(\mathrm{F} = \mathrm{F}(\mathrm{X})\) 借助同态 \(g\) 嵌入。滤过 \((\Gamma_n)\) 的 \(\Gamma\) 诱导出滤过 \((\mathbb{F}_n)\)，该滤过定义在 \(\mathbf{F}\) 上，这是一个受限滤过。李 \(p\)-代数 \(\operatorname{gr}(\mathbf{F})\) 可识别为一个李 \(p\)-子代数 \(\mathrm{A}(\mathbf{X})\)，其中含有 \(\mathbf{X}\)。推出 \(\operatorname{gr}(\mathbf{F})\) 含有自由李 \(p\)-代数 \(\mathrm{L}(\mathbf{X}, p)\)，参见~§ 3，习题 4(e)。
\item
  令 \(\mathbf{H}\) 是关于 \(\mathbf{X}\) 的 Hall 集；对每个整数 \(i\)，令 \(\mathbf{H}_i\) 为 \(\mathbf{H}\) 中长度为 \(i\) 的元素所成的集合。令 \(w \in \mathbf{F}\)，并令 \(\alpha_i \in \mathbf{Z}^{(\mathbf{H}_i)}\) 满足
\end{enumerate}

\[
w \equiv \prod_ {i = 1} ^ {n} \prod_ {m \in H_i} \phi (m) ^ {\alpha_ {i} (m)} \mod C ^ {n + 1} F
\]

对所有 \(n\) 均成立，参见~第 4 号。对每个 \(m \in \mathbf{H}\)，令 \(l(m)\) 表示 \(m\) 的长度，令 \(q(m, w)\) 表示 \(p\) 的最大次幂，使其整除 \(\alpha_{1}(m)\)，其中 \(i = l(w)\)（若 \(\alpha_{1}(m) = 0\)，约定 \(q(m, w) = +\infty\)）。令 \(\mathrm{N} = \inf_{m \in \mathbb{H}} (l(m).q(m, w))\)。证明 \(\phi(m)^{\alpha_{1}(m)}\) 属于 \(\mathrm{F}_{\mathrm{N}}\)，并且甚至属于 \(\mathrm{F}_{\mathrm{N+1}}\)（当 \(l(w).q(m, w) > \mathrm{N}\) 时）。证明若 \(l(w).q(m, w) = \mathrm{N}\)，则 \(\phi(m)^{\alpha_{1}(m)}\) 在 \(\mathrm{gr}_{\mathrm{N}}(\hat{\mathrm{A}}(\mathbf{X})) = \mathrm{A}^{\mathrm{N}}(\mathbf{X})\) 中的像是 \(\overline{m}^{q(m, w)}\) 的非零倍数，其中 \(\overline{m}\) 是 \(\mathrm{L}(\mathbf{X})\) 中由 \(m\) 定义的元素（\(\S 2\)，第 11 号）。现在这些元素线性无关（\(\S 3\)，习题 4）；推出 \(w\) 在 \(\mathrm{gr}_{\mathrm{N}}(\mathbf{F})\) 中的像是 \(\neq 0\) 且属于 \(\mathrm{L}(\mathbf{X}, p)\)，从而 \(\mathrm{gr}(\mathbf{F}) = \mathrm{L}(\mathbf{X}, p)\)。

\begin{enumerate}
\def\labelenumi{\arabic{enumi}.}
\setcounter{enumi}{11}
\tightlist
\item
  令 g 是特征为 p \textgreater{} 0 的域 k 上的李代数。
\end{enumerate}

\begin{enumerate}
\def\labelenumi{(\alph{enumi})}
\tightlist
\item
  令 \(h\) 为满足 \(\leqslant p - 1\) 的整数。\(\mathfrak{g}\) 称为满足 Engel 的第 \(h\) 条件，是指 \((\operatorname{ad} x)^h(y) = 0\) 对每个有序对 \((x, y)\)（其元素属于 \(\mathfrak{g}\)）都成立。证明这等价于
\end{enumerate}

\[
\sum_ {\sigma \in \mathfrak {S} _ {h}} \operatorname{ad} (x _ {\sigma (1)}) \circ \dots \circ \operatorname{ad} (x _ {\sigma (h)}) = 0
\]

当 \(x_1, \ldots, x_h\) 属于 \(\mathfrak{g}\) 时成立（将公式 \((\operatorname{ad} x)^h = 0\) 应用于 \(x_i\) 的线性组合）。


\begin{enumerate}
\def\labelenumi{(\alph{enumi})}
\setcounter{enumi}{1}
\item
  令 \(x, y\) 为 \(\mathfrak{g}\) 中的元素，满足 \(\Lambda_p(ax, by) = 0\)，其中 \(a, b\) 属于 \(k\)。证明若 \(\Lambda_p^{r,s}\) 是 \(\Lambda_p\) 的双齐次分量，且双次数为 \((r, s)\)，则 \(\Lambda_p^{r,s}(x, y) = 0\) 对每个有序对 \((r, s)\) 都成立。推出（取 \(r = p - 1\)、\(s = 1\)）\((\text{ad } x)^{p-1}(y) = 0\)。特别地，若 \(\Lambda_p(x, y) = 0\)，其中 \(x, y\) 属于 \(\mathfrak{g}\)，则李代数 \(\mathfrak{g}\) 满足 Engel 的第 \((p - 1)\) 条件。
\item
  令 \(\mathfrak{c}\) 为 \(\mathfrak{g}\) 的中心。假设 \((\operatorname{ad} x)^p = 0\) 对所有 \(x \in \mathfrak{g}\) 都成立。证明 \(\mathfrak{g} / \mathfrak{c}\) 满足 Engel 的第 \((p - 1)\) 条件。（证明 \(\Lambda_p(\operatorname{ad} x, \operatorname{ad} y) = 0\) 对 \(x, y\) 属于 \(\mathfrak{g}\) 成立，并将 \((b)\) 应用于李代数 \(\operatorname{ad} \mathfrak{g} = \mathfrak{g} / c\)。）
\item
  令 G 是满足 \(x^{p}=1\) 对所有 \(x\in G\) 成立的群，令 \((\mathrm{G}_{\alpha})\) 是 G 的实中心滤过。滤过 \((\mathrm{G}_{\alpha})\) 是受限的（习题 10），而 \(\mathrm{gr}(\mathrm{G})\) 是 p-映射为零的李 p-代数。推出 \(\mathrm{gr}(\mathrm{G})\) 满足 Engel 的第 \((p-1)\) 条件。†
\end{enumerate}

\exercisesfor{6}

\begin{enumerate}
\def\labelenumi{\arabic{enumi}.}
\tightlist
\item
  证明 \(\exp (\mathrm{U}).\exp (\mathrm{V}).\exp (-\mathrm{U}) = \exp \left(\sum_{n = 0}^{\infty}\frac{1}{n!} (\operatorname {ad}\mathrm{U})^{n}(\mathrm{V})\right)\)。推出
\end{enumerate}

\[
\mathrm{H} (\mathrm{U}, \mathrm{H} (\mathrm{V}, - \mathrm{U})) = \sum_ {n = 0} ^ {\infty} \frac {1}{n !} (\operatorname{ad} \mathrm{U}) ^ {n} (\mathrm{V}) = (\exp (\operatorname{ad} \mathrm{U})) (\mathrm{V}).
\]

\begin{enumerate}
\def\labelenumi{\arabic{enumi}.}
\setcounter{enumi}{1}
\item
  证明 \(\mathbf{H}(\mathbf{U},\mathbf{V}) = \sum_{n = 0}^{\infty}\frac{1}{n!} (\operatorname {ad}\mathbf{U})^{n}(\mathbf{H}(\mathbf{V},\mathbf{U}))\)。（应用习题 1，注意 \(\mathbf{H}(\mathbf{U},\mathbf{H}(\mathbf{H}(\mathbf{V},\mathbf{U}), - \mathbf{U})) = \mathbf{H}(\mathbf{U},\mathbf{V}).\)）
\item
  令 \(\mathbf{M}(\mathbf{U},\mathbf{V}) = \mathbf{H}(-\mathbf{U},\mathbf{U} + \mathbf{V})\)。于是
\end{enumerate}

\[
\exp (\mathrm{M} (\mathrm{U}, \mathrm{V})) = \exp (- \mathrm{U}). \exp (\mathrm{U} + \mathrm{V}).
\]

令 \(M_{1}(U, V)\)（相应地，\(H_{1}(U, V)\)）表示级数 M（相应地，H）中关于 V 的次数为 1 的双齐次项之和。

\begin{enumerate}
\def\labelenumi{(\alph{enumi})}
\tightlist
\item
  证明
\end{enumerate}

\[
\mathrm{M} _ {1} (\mathrm{U}, \mathrm{V}) = \sum_ {n = 0} ^ {\infty} \frac {(- 1) ^ {n}}{(n + 1) !} (\mathrm{adU}) ^ {n} (\mathrm{V}) = (f (\mathrm{adU})) (\mathrm{V})
\]

其中 \(f(\mathrm{T}) = (1 - e^{-\mathrm{T}}) / \mathrm{T}\)。

（论证与命题 5 的证明相同。也可以利用恒等式将其归结为命题 5：

\[
\exp (\mathrm{U}). \exp (\mathrm{M} (\mathrm{U}, \mathrm{V})). \exp (- \mathrm{U}) = \exp (\mathrm{H} (\mathrm{U} + \mathrm{V}, - \mathrm{U})),
\]

并结合习题 1。）

\begin{enumerate}
\def\labelenumi{(\alph{enumi})}
\setcounter{enumi}{1}
\tightlist
\item
  证明 \(\mathrm{H}(\mathrm{U},\mathrm{M}(\mathrm{U},\mathrm{V})) = \mathrm{U} + \mathrm{V}\)。推出
\end{enumerate}

\[
\mathrm{H} _ {1} (\mathrm{U}, \mathrm{M} _ {1} (\mathrm{U}, \mathrm{V})) = \mathrm{V}.
\]

\begin{enumerate}
\def\labelenumi{(\alph{enumi})}
\setcounter{enumi}{2}
\tightlist
\item
  令 \(g(\mathrm{T}) = 1 / f(\mathrm{T}) = 1 + \mathrm{T} / 2 + \sum_{n=1}^{\infty} \frac{1}{(2n)!} b_{2n}\mathrm{T}^{2n}\)，其中 \(b_{2n}\) 是
\end{enumerate}

伯努利数（《实变函数》，第六章，§ 1，第 4 号）。证明

\[
\mathrm{H} _ {1} (\mathrm{U}, \mathrm{V}) = (g (\operatorname{ad} \mathrm{U})) (\mathrm{V}) = \mathrm{V} + \frac {1}{2} [ \mathrm{U}, \mathrm{V} ] + \sum_ {n = 1} ^ {\infty} \frac {1}{(2 n) !} b _ {2 n} (\operatorname{ad} \mathrm{U}) ^ {2 n} (\mathrm{V}).
\]

写出 \(H_{1}(U, V)\) 展开的前几项：

\[
\begin{array}{r l} \mathrm{H} _ {1} (\mathrm{U}, \mathrm{V}) & = \mathrm{V} + \frac {1}{2} (\mathrm{adU}) (\mathrm{V}) + \frac {1}{12} (\mathrm{adU}) ^ {2} (\mathrm{V}) - \frac {1}{720} (\mathrm{adU}) ^ {4} (\mathrm{V}) \\ & + \frac {1}{30240} (\mathrm{adU}) ^ {6} (\mathrm{V}) - \dots \end{array}
\]

\begin{enumerate}
\def\labelenumi{(\alph{enumi})}
\setcounter{enumi}{3}
\tightlist
\item
  证明 H(U, V) 中关于 U 的次数为 1 的双齐次项之和为级数：
\end{enumerate}

\[
\mathrm{U} + \frac {1}{2} [ \mathrm{U}, \mathrm{V} ] + \sum_ {n = 1} ^ {\infty} \frac {1}{(2 n) !} b _ {2 n} (\mathrm{adV}) ^ {2 n} (\mathrm{U}).
\]

（将 (c) 应用于 H(V, U)，并使用习题 2 将 H(V, U) 转换为 H(U, V)。）(e) 令 W 为另一个不定元，令 X(U, V, W) 为 H(U, V + W) 中关于 W 的次数为 1 的三齐次项之和。证明

\[
\mathrm{X} (\mathrm{U}, \mathrm{H} _ {1} (\mathrm{V}, \mathrm{W})) = \mathrm{H} _ {1} (\mathrm{H} (\mathrm{U}, \mathrm{V}), \mathrm{W}).
\]

（使用恒等式 \(\mathrm{H}(\mathrm{U},\mathrm{H}(\mathrm{V},\mathrm{W})) = \mathrm{H}(\mathrm{H}(\mathrm{U},\mathrm{V}),\mathrm{W}).\)）推出

\[
\mathrm{X} (\mathrm{U}, \mathrm{V}, \mathrm{W}) = (g (\operatorname{ad} \mathrm{H} (\mathrm{U}, \mathrm{V})) \circ f (\operatorname{ad} \mathrm{V})) (\mathrm{W}).
\]

¶ 4. 令 X 是有限集，令 \(\hat{\mathbf{L}} = \hat{\mathbf{L}}(\mathbf{X})\)。在 \(\hat{\mathbf{L}}\) 上赋予 Hausdorff 复合律 \((a, b) \mapsto a \uplus b\)，并简记为 \(a.b\)。若 \(t \in \mathbf{K}\)、\(a \in \hat{\mathbf{L}}\)，则记 \(a^t = t a\)。

\begin{enumerate}
\def\labelenumi{(\alph{enumi})}
\tightlist
\item
  令 \(\mathcal{H}\) 是关于 X 的 Hall 集。若 \(m \in \mathcal{H}\)，令 \(m_{\mathrm{L}}\) 表示李代数 \(\mathrm{L}(\mathrm{X})\) 中相应的元素（§2，第 11 号），令 \(m_{\mathcal{H}}\) 表示群 \(\hat{\mathrm{L}}\) 的基本交换子，它由 \(m\) 定义（§5，第 4 号，注）。若 \(m\) 的长度为 \(l\)，证明 \(m_{\mathcal{H}} = m_{\mathrm{L}} + \mu\)，其中 \(\mu\) 的滤过次数 \(\geqslant l + 1\)（对 \(l\) 作归纳）。推出每个元素 \(w\) 都属于群 \(\hat{\mathrm{L}}\)，并且能唯一地写成
\end{enumerate}

\[
w = \prod_ {m \in \mathcal {H}} m _ {\mathcal {H}} ^ {\alpha (m)},
\]


其中 \(\alpha(m) \in \mathbf{K}\)，且该乘积在拓扑群 \(\hat{\mathbf{L}}\) 中收敛。

\begin{enumerate}
\def\labelenumi{(\alph{enumi})}
\setcounter{enumi}{1}
\item
  令 \(\mathbf{P}\) 为素数集合，令 \(c\) 为满足 \(\geqslant 1\) 的整数。令 \(\mathbf{K} = \mathbf{Q}\)。证明群 \(\hat{\mathbf{L}} / \hat{\mathbf{L}}_c\) 可识别为 \(\mathbf{P}\)-包络 \(\mathrm{F}_{\mathbf{P}}^{c}\)，其所包络的群为 \(\mathrm{F}^c = \mathrm{F}(\mathrm{X}) / \mathrm{C}^c\mathrm{F}(\mathrm{X})\)，参见~§ 5，习题 6。
\item
  取 X 为含两个元素 \(\{U, V\}\) 的集合。证明存在两个由有理数组成的族 \(\alpha(m)\)、\(\beta(m)\)，使得
\end{enumerate}

\[
\mathrm{U} + \mathrm{V} = \prod_ {m \in \mathcal {H}} m _ {\mathcal {H}} ^ {\alpha (m)}
\]

\[
[ \mathrm{U}, \mathrm{V} ] = \prod_ {m \in \mathcal {H}} m _ {\mathcal {H}} ^ {\beta (m)}.
\]

例如

\[
\begin{array}{c} \mathrm{U} + \mathrm{V} = \mathrm{U}. \mathrm{V}. (\mathrm{U}, \mathrm{V}) ^ {- \frac {1}{2}} \dots \\ {} [ \mathrm{U}, \mathrm{V} ] = (\mathrm{U}, \mathrm{V}) \dots \end{array}
\]

\begin{enumerate}
\def\labelenumi{(\alph{enumi})}
\setcounter{enumi}{3}
\tightlist
\item
  令 g 是带有 Hausdorff 群律的幂零李 Q-代数，令 u、v 属于 g。证明
\end{enumerate}

\[
u + v = \prod_ {m \in \mathcal {H}} m _ {\mathcal {H}} (u, v) ^ {\alpha (m)}
\]

\[
[ u, v ] = \prod_ {m \in \mathcal {H}} m _ {\mathcal {H}} (u, v) ^ {\beta (m)},
\]

其中 \(\alpha\) 和 \(\beta\) 是上面定义的有理数族（“Hausdorff 公式的逆”）。（使用连续同态 \(\phi\colon\hat{L}\to\mathfrak{g}\)，满足 \(\phi(U)=u\)、\(\phi(V)=v\)，并注意到它对于 Hausdorff 群律也是同态。）

\begin{enumerate}
\def\labelenumi{(\alph{enumi})}
\setcounter{enumi}{4}
\tightlist
\item
  令 G 是无挠可除幂零群（即 P-无挠且 P-可除）。若 \(u \in G\)、\(t \in Q\)，则 \(u^{t}\) 已定义（\(\S 4\)，习题 16）。若 \(u \in G\)、\(v \in G\)，则用上面的 (d) 中的公式定义 \(u + v\) 和 {[}u, v{]}。证明这样 G 具有幂零李 Q-代数结构，且相应的 Hausdorff 律就是 G 上给定的群律（当 G 是群 \(F_{P}^{c}\) 时验证这些断言，参见~(b)），再利用从 \(F_{P}^{r}\) 到 G 的同态将其推广到一般情形。
\end{enumerate}

令 \(f\) 是从 \(G\) 到无挠可除幂零群 \(G'\) 的映射。证明 \(f\) 是（群）同态，当且仅当它是李代数同态。（因此，Hausdorff 律在幂零李 \(\mathbf{Q}\)-代数范畴与无挠可除幂零群范畴之间定义了一个同构。）

\begin{enumerate}
\def\labelenumi{\arabic{enumi}.}
\setcounter{enumi}{4}
\tightlist
\item
  令 \(\mathfrak{g}\) 是赋予 Hausdorff 群律的幂零李 \(\mathbf{Q}\)-代数，将该群律记为 \((x,y)\mapsto x.y\)。
\end{enumerate}

\begin{enumerate}
\def\labelenumi{(\alph{enumi})}
\tightlist
\item
  证明若 \(x \in \mathfrak{g}\)、\(y \in \mathfrak{g}\)，则
\end{enumerate}

\[
x. y. x ^ {- 1} = \sum_ {n = 0} ^ {\infty} \frac {1}{n !} (\operatorname{ad} x) ^ {n} (y).
\]

（使用习题 2。）

\begin{enumerate}
\def\labelenumi{(\alph{enumi})}
\setcounter{enumi}{1}
\item
  令 \(\mathfrak{h}\) 是 \(\mathfrak{g}\) 的子集。证明 \(\mathfrak{h}\) 是 \(\mathfrak{g}\) 的李子代数（相应地，理想），当且仅当它是群 \(\mathfrak{g}\) 的饱和子群（相应地，饱和正规子群）。（使用习题 4(d) 的公式，将群律转换为李代数律。）
\item
  令 \(\mathfrak{h}\) 是群 \(\mathfrak{g}\) 的子群。证明 \(\mathbf{P}\)-饱和的 \(\mathfrak{h}\) 是 \(\mathbf{Q}.\mathfrak{h}\)。
\item
  令 \(\mathfrak{h}\) 是 \(\mathfrak{g}\) 的李子代数。证明 \(\mathfrak{h}\) 在群 \(\mathfrak{g}\) 中的中心化子（相应地，正规化子）是所有 \(x \in \mathfrak{g}\) 中满足 \((\operatorname{ad} x)(\mathfrak{h}) = 0\)（相应地，\((\operatorname{ad} x)(\mathfrak{h}) \subset \mathfrak{h}\)）者所成的集合。
\item
  证明李代数 g 的下中心列与群 g 的下中心列重合，并且相关的分次李代数 gr(g) 从“群”的观点看与从“李代数”的观点看相同。
\end{enumerate}

¶ 6. 令 G 是有限生成的无挠幂零群。证明当 n 充分大时，G 可嵌入 Z 上的 n 阶严格下三角群中。（令 \((i, \overline{G})\) 为 G 的一个 P-包络，参见~§ 5，习题 6；赋予其典范李代数结构后（习题 4），\(\overline{G}\) 是有限维幂零李 Q-代数。将 Ado 定理（第一章，§ 7）应用于 \(\overline{G}\)，推出 \(\overline{G}\) 可嵌入 Q 上的严格下三角群。再通过适当的整对角矩阵共轭，将其从 Q 转到 Z。）

\exercisesfor{7}

\begin{enumerate}
\def\labelenumi{\arabic{enumi}.}
\tightlist
\item
  令 \(\mathfrak{g}\) 是 \(\mathbf{K}\) 上满足下式的完备赋范李代数：
\end{enumerate}

\[
\| [ x, y ] \| \leqslant \| x \| \| y \|
\]

当 \(x, y\) 属于 \(\mathfrak{g}\) 时成立。令 \(\Theta\) 表示所有 \(x \in \mathfrak{g}\) 中满足 \(\| x \| < \frac{1}{3} \log \frac{3}{2}\) 者所成的集合。对于 \(x, y\) 属于 \(\Theta\)，有 \(h(x, y) \in \Theta\)，参见~第 2 号。

\begin{enumerate}
\def\labelenumi{(\alph{enumi})}
\item
  记 \(f(\mathrm{T}) = (1 - \varepsilon^{-\mathrm{T}}) / \mathrm{T}\)，\(g(\mathrm{T}) = 1 / f(\mathrm{T})\)。级数 \(f\) 在整个复平面上收敛，而级数 \(g\) 在半径为 \(2\pi\) 的开圆盘上收敛（参见~《实变函数》，第六章，§2，第 3 号）。推出若 \(z \in \Theta\)，则 \(f(\mathrm{ad}z)\) 和 \(g(\mathrm{ad}z)\) 有定义，且它们是 \(\mathcal{L}(\mathfrak{g};\mathfrak{g})\) 的元素，并互为逆元。
\item
  令 \(\mathrm{D}_2h(x,y)\) 为 \(h\) 在点 \((x,y)\)（该点属于 \(\Theta \times \Theta\)）处的第二个偏导数（《微分与解析流形》，R，1.6.2）。它是 \(\mathcal{L}(\mathfrak{g};\mathfrak{g})\) 的元素。证明
\end{enumerate}

\[
\mathrm{D} _ {2} h (x, y) = g (\operatorname{ad} h (x, y)) \circ f (\operatorname{ad} y).
\]

（使用§6 的习题 3(e)，证明当 \(x\)、\(y\) 充分接近零时该公式成立，再通过解析延拓将其推广到一般情形。）证明公式

\[
f (\operatorname{ad} h (x, y)) \circ \mathrm{D} _ {2} h (x, y) = f (\operatorname{ad} y)
\]

在形式幂级数 \(\tilde{H}\) 的绝对收敛域的每一点都成立（参见~命题 1）。

\exercisesfor{8}

设 \(p\) 是域 \(K\) 的剩余特征，且 \(>0\)。记 \(\mathfrak{o}_{K}\) 为赋值 \(v\) 在域 \(K\) 上的环。

\begin{enumerate}
\def\labelenumi{\arabic{enumi}.}
\tightlist
\item
  \begin{enumerate}
  \def\labelenumii{(\alph{enumii})}
  \tightlist
  \item
    若元素 \(\pi\) 属于 \(K\) 且满足 \(v(\pi) = \theta\) 以及 \(\pi^{p-1} / p \equiv -1 \pmod{\pi o_K}\)，则称其为可容许元素。
  \end{enumerate}
\end{enumerate}

（注意 \(\pi^{p-1}/p\in\mathfrak{o}_{K}\)，因为 \(v(\pi)=\theta.\)）

证明存在满足本段条件且含有可容许元素的域 K（向 \(Q_{p}\) 添加一个 \((p - 1)\) 次根 \(-p\)）。

\begin{enumerate}
\def\labelenumi{(\alph{enumi})}
\setcounter{enumi}{1}
\tightlist
\item
  令 \(\pi\) 是 K 中的可容许元素。考虑下列系数属于 K 的形式幂级数：
\end{enumerate}

\[
e _ {\pi} (\mathrm{U}) = \frac {1}{\pi} e (\pi \mathrm{U}) = \sum_ {n = 1} ^ {\infty} \pi^ {n - 1} \mathrm{U} ^ {n} / n!,
\]

\[
l _ {\pi} (\mathrm{U}) = \frac {1}{\pi} l (\pi \mathrm{U}) = \sum_ {n = 1} ^ {\infty} (- \pi) ^ {n - 1} \mathrm{U} ^ {n} / n.
\]

证明这些级数互为逆级数，且其系数属于 \(\mathfrak{o}_{K}\)（使用引理 2）。

\begin{enumerate}
\def\labelenumi{(\alph{enumi})}
\setcounter{enumi}{2}
\tightlist
\item
  令 \(\{\mathrm{U},\mathrm{V}\}\) 是含两个元素的集合，令 \(\mathrm{H}(\mathrm{U},\mathrm{V})\) 为 Hausdorff 级数。记
\end{enumerate}

\[
\mathrm{H} _ {\pi} (\mathrm{U}, \mathrm{V}) = \frac {1}{\pi} \mathrm{H} (\pi \mathrm{U}, \pi \mathrm{V}).
\]

于是 \(H_{\pi}(U, V) \in \hat{L}_{K}(\{U, V\})\)。

证明

\[
\mathrm{H} _ {\pi} (\mathrm{U}, \mathrm{V}) = l _ {\pi} (e _ {\pi} (\mathrm{U}) + e _ {\pi} (\mathrm{V}) + \pi e _ {\pi} (\mathrm{U}). e _ {\pi} (\mathrm{V})).
\]

推出 \(H_{\pi} \in \hat{L}_{\mathfrak{o}_{K}}(\{\mathrm{U}, \mathrm{V}\})\)，即 \(H_{\pi}\) 的系数属于 \(o_{K}\)；从而得到命题 1 的另一种证明。

\begin{enumerate}
\def\labelenumi{(\alph{enumi})}
\setcounter{enumi}{3}
\tightlist
\item
  令 \(\tilde{e}_{\pi}, \tilde{l}_{\pi}\) 和 \(\tilde{\mathbf{H}}_{\pi}\) 表示将 \(e_{\pi}, l_{\pi}\) 和 \(\mathbf{H}_{\pi}\) 模 \(\pi \mathfrak{o}_{K}\) 化简所得的级数；它们的系数属于环 \(\mathfrak{o}_{K} / \pi \mathfrak{o}_{K}\)。
\end{enumerate}

证明 \(\tilde{l}_{\pi}(\mathrm{U}) = \mathrm{U} - \mathrm{U}^{p}\)（注意 \(v_{p}(n) < (n - 1)\theta\) 对 \(n \neq 1, p\) 成立）。推出

\[
\tilde {e} _ {\pi} (\mathrm{U}) = \mathrm{U} + \mathrm{U} ^ {p} + \dots + \mathrm{U} ^ {p ^ {n}} + \dots
\]

并且

\[
\tilde {\mathrm{H}} _ {\pi} (\mathrm{U}, \mathrm{V}) = \tilde {e} _ {\pi} (\mathrm{U}) + \tilde {e} _ {\pi} (\mathrm{V}) - (\tilde {e} _ {\pi} (\mathrm{U}) + \tilde {e} _ {\pi} (\mathrm{V})) ^ {p},
\]

或者也可写为

\[
\tilde {\mathrm{H}} _ {\pi} (\mathrm{U}, \mathrm{V}) = \mathrm{U} + \mathrm{V} - \Lambda_ {p} \left(\sum_ {n = 0} ^ {\infty} \mathrm{U} ^ {p ^ {n}}, \sum_ {n = 0} ^ {\infty} \mathrm{V} ^ {p ^ {n}}\right),
\]

其中 \(\Lambda_{p}\) 由下式定义：

\[
\Lambda_ {p} (\mathrm{U}, \mathrm{V}) = (\mathrm{U} + \mathrm{V}) ^ {p} - \mathrm{U} ^ {p} - \mathrm{V} ^ {p}, \quad \text { cf.   Chapter   I,   } \S 1, \text {   Exercise   19. }
\]

特别地，\(\tilde{\mathrm{H}}_{\pi}(\mathrm{U},\mathrm{V})\) 属于 \(\hat{\mathrm{L}}_{\mathbf{F}_{p}}(\{\mathrm{U},\mathrm{V}\})\)，且其 \(p\) 次齐次分量为 \(-\Lambda_p(\mathrm{U},\mathrm{V})\)。于是

\[
\begin{array}{c} \tilde {\mathrm{H}} _ {\pi} (\mathrm{U}, \mathrm{V}) = \tilde {\mathrm{H}} _ {\pi} (\mathrm{V}, \mathrm{U}) \\ \tilde {\mathrm{H}} _ {\pi} (\mathrm{U}, \tilde {\mathrm{H}} _ {\pi} (\mathrm{V}, \mathrm{W})) = \tilde {\mathrm{H}} _ {\pi} (\tilde {\mathrm{H}} _ {\pi} (\mathrm{U}, \mathrm{V}), \mathrm{W}) \quad \text { in } \hat {\mathrm{L}} _ {\mathbf {F} _ {p}} (\{\mathrm{U}, \mathrm{V}, \mathrm{W} \}), \end{array}
\]

其中 W 是第三个不定元。

\begin{enumerate}
\def\labelenumi{(\alph{enumi})}
\setcounter{enumi}{4}
\tightlist
\item
  令 C(U, V) = H(-U, H(-V, H(U, V)))（“Hausdorff 交换子”）。于是 exp(C(U, V)) = exp(-U) exp(-V) exp(U) exp(V)。记
\end{enumerate}

\[
\mathrm{C} _ {\pi} (\mathrm{U}, \mathrm{V}) = \frac {1}{\pi} \mathrm{C} (\pi \mathrm{U}, \pi \mathrm{V}).
\]

证明 \(C_{\pi}(U,V)\) 的系数属于 \(\pi o_{K}\)（使用 \(\tilde{H}_{\pi}(U,V)=\tilde{H}_{\pi}(V,U)\) 这一事实）。†

\begin{enumerate}
\def\labelenumi{\arabic{enumi}.}
\setcounter{enumi}{1}
\tightlist
\item
  我们知道（§ 6，习题 3），若 \(n \geqslant 2\)，则 H(U, V) 的双次数为 \((n, 1)\) 的分量是 \(\frac{1}{n!} b_n(\mathrm{ad}\mathrm{U})^n (\mathrm{V})\)，其中 \(b_n\) 是第 \(n\) 个伯努利数。
\end{enumerate}

由此和命题 1 推出不等式

\[
v _ {p} (b _ {n} / n!) \leqslant n / (p - 1).
\]

利用克劳森—冯·施陶特定理（《实变函数》，第六章，§2，习题 6）重新得到这一结果，并证明等号成立当且仅当 \(S(n) = p - 1\)。

\begin{enumerate}
\def\labelenumi{\arabic{enumi}.}
\setcounter{enumi}{2}
\tightlist
\item
  假设 K 含有本原 \(p\) 次单位根 \(w\)。记 \(\pi = w - 1\)。利用公式
\end{enumerate}

\[
w ^ {i} - 1 = \pi (1 + w + \dots + w ^ {i - 1})
\]

证明 \(v(w^{i} - 1) = v(\pi)\) 对 \(1 \leqslant i \leqslant p - 1\) 成立，并且

\[
\frac {w ^ {i} - 1}{\pi} \equiv i \pmod {\pi}.
\]

利用公式 \(p = \prod_{i=1}^{p-1} (w^i - 1)\) 推出 \(\pi\) 是可容许元素（习题 1）。

¶ 4. 令 c 是满足 \(\geqslant1\) 的整数，令 P 是包含所有 \(\leqslant c\) 的素数的素数集合。令 \(Z_{P}=S_{P}^{-1}Z\)（§ 4，习题 16）。证明 Hausdorff 级数 H(U, V) 中次数 \(\leqslant c\) 的项属于 \(L_{Z_{P}}(\{U, V\})\)。推出，若 g 是幂零李 \(Z_{P}\)-代数，且类为 \(\leqslant c\)，则复合律 \((u,v)\mapsto\mathrm{H}(u,v)\) 使 g 成为类 \(\leqslant c\) 的 P-无挠、P-可除幂零群。证明反过来每个具有这些性质的群都可由此过程得到。（使用“Hausdorff 公式的逆”，参见~§6，习题 4，并证明其中出现的指数 \(\alpha(m)\)、\(\beta(m)\) 属于 \(Z_{P}\)，当 \(l(m)\leqslant c\) 时。）

特别地，每个阶为 \(p^{n}\) 且类 \textless p 的 p-群，都可以通过 Hausdorff 律由一个类 \textless p 且含有 \(p^{n}\) 个元素的幂零李 Z-代数得到。（取 P 为异于 p 的素数集合。）

\exercisesforappendix

\begin{enumerate}
\def\labelenumi{\arabic{enumi}.}
\tightlist
\item
  令 \(\Phi_n\) 为指标为 \(n\) 的圆分多项式（《代数》，第五章，§ 11，第 2 号）。利用公式
\end{enumerate}

\[
\mathbf {X} ^ {n} - 1 = \prod_ {d | n} \Phi_ {d} (\mathbf {X}),
\]

证明

\[
\Phi_ {n} (\mathrm{X}) = \prod_ {d | n} \left(\mathrm{X} ^ {n / d} - 1\right) ^ {\mu (d)}.
\]

（将 Möbius 反演公式应用于域 \(\mathbf{Q}(\mathbf{X})\) 的乘法群。）

\begin{enumerate}
\def\labelenumi{\arabic{enumi}.}
\setcounter{enumi}{1}
\tightlist
\item
  令 D 为幺半群 \(\mathbf{N}^*\) 的全代数（《代数》，第三章，§ 2，第 10 号）。若 \(n \in \mathbf{N}^*\)，令 \(n^\omega\) 表示它在 D 中的像，于是 \(1^\omega = 1\)，且 \((nm)^\omega = n^\omega m^\omega\) 对 \(n\)、\(m\) 属于 \(\mathbf{N}^*\) 成立。D 中每个元素 \(f\) 都能唯一地写成级数 \(\sum_{i=1}^{\infty}\)
\end{enumerate}

\(\sum_{n=1}^{\infty} a_n n^\omega\)，其中 \(a_n \in \mathbf{K}\)。

\begin{enumerate}
\def\labelenumi{(\alph{enumi})}
\item
  令 \(f = \sum_{n=1}^{\infty} a_n n^\omega\) 是 D 中的元素。证明 \(f\) 在 D 中可逆，当且仅当 \(a_1\) 在 K 中可逆。特别地，若 K 是局部环，则 D 也是局部环。
\item
  记 \(\zeta = \sum_{n=1}^{\infty} n^{\omega}\)，\(\mu = \sum_{n=1}^{\infty} \mu(n)n^{\omega}\)。证明 \(\zeta\) 与 \(\mu\) 互为逆元。推出若 \(s = \sum_{n} s(n)n^{\omega}\) 和 \(t = \sum_{n} t(n)n^{\omega}\) 是 D 中的两个元素，则关系 \(s = \zeta\) . \(t\) 与 \(t = \mu\) . \(s\) 等价（Möbius 反演公式的一个变体）。
\item
  令 \(\mathbf{P}\) 为素数集合。证明族 \((1 - p^{\omega})\) 对 \(p\in \mathbf{P}\) 可在 \(\mathbf{D}\) 中相乘，并且
\end{enumerate}

\[
\mu = \sum_ {p \in P} ^ {\infty} (1 - p ^ {\omega}) \quad \text { and } \quad \zeta = \prod_ {p \in P} \frac {1}{1 - p ^ {\omega}}.
\]

\chapter{李群}

本章中，K 表示以下三者之一：实数的带赋值域 R、复数的带赋值域 C，或非离散完备超度量交换域。我们假定从 § 4 起 K 的特征为 0，§ 6 中 K = R 或 C，§ 7 中 K 为超度量域。除非另有说明，所考虑的所有流形、代数和向量空间都定义在 K 上。回忆一下，当我们说一个 \(C^{r}\) 类流形时，r ∈ N\_K，即若 K ≠ R 则 r = ω，而若 K = R 则 l ≤ r ≤ ω。

关于范数、可赋范空间和赋范空间的约定采用《微分与解析流形》R 中的约定。

回忆一下，K 上的可赋范代数是 K 上的一个（不一定结合的）代数 A，带有满足下列性质的拓扑 T：

\begin{enumerate}
\def\labelenumi{(\arabic{enumi})}
\item
  \(\mathcal{T}\) 可由范数定义；
\item
  映射 \((x,y)\mapsto xy\) 从 \(\mathbf{A}\times \mathbf{A}\) 到 \(\mathbf{A}\) 是连续的。
\end{enumerate}

A 的双连续自同构群记为 \(\mathrm{Aut(A)}\)。每个有限维 \(K\) 上的代数都是带有典范拓扑的可赋范代数。\(K\) 上的赋范代数是带有范数的 \(K\) 上代数 A，满足范数不等式 \(\| xy\| \leqslant \| x\|\| y\|\) 对 A 中所有 \(x,y\) 都成立；赋予由该范数定义的拓扑后，代数 A 是可赋范代数。若 A 是可赋范代数，则 A 上存在一个定义其拓扑并使其成为赋范代数的范数。

若 G 是群，则 \(e_{G}\)，或简写为 e，表示 G 的单位元。对 \(g \in G\)，\(\gamma(g)\)、\(\tilde{\delta}(g)\) 和 \(\operatorname{Int}(g)\) 分别表示 G 到 G 的映射 \(g' \mapsto gg'\)、\(g' \mapsto g'g^{-1}\) 和 \(g' \mapsto gg'g^{-1}\)。若 f 是从 G 到集合 E 的映射，则 f 表示从 G 到 E 的映射 \(g \mapsto f(g^{-1})\)。

\section{李群}
